\documentclass[11pt]{article}

\usepackage[margin=1in]{geometry}
\usepackage{bm}
\usepackage{amsmath}
\usepackage{amsfonts}
\usepackage{amssymb}
\usepackage{amsthm}
\usepackage{mathtools}
\usepackage{enumitem,microtype}
\usepackage[hidelinks]{hyperref}
\usepackage[nameinlink,capitalise,noabbrev]{cleveref}
\usepackage{aliascnt}

\newtheorem{theorem}{Theorem}[section]
\newaliascnt{proposition}{theorem}
\newtheorem{proposition}[proposition]{Proposition}
\aliascntresetthe{proposition}
\newaliascnt{lemma}{theorem}
\newtheorem{lemma}[lemma]{Lemma}
\aliascntresetthe{lemma}
\newaliascnt{corollary}{theorem}
\newtheorem{corollary}[corollary]{Corollary}
\aliascntresetthe{corollary}
\newaliascnt{remark}{theorem}
\newtheorem{remark}[remark]{Remark}
\aliascntresetthe{remark}
\theoremstyle{definition}
\newaliascnt{definition}{theorem}
\newtheorem{definition}[definition]{Definition}
\aliascntresetthe{definition}

\crefname{theorem}{theorem}{theorems}
\Crefname{theorem}{Theorem}{Theorems}
\crefname{proposition}{proposition}{propositions}
\Crefname{proposition}{Proposition}{Propositions}
\crefname{lemma}{lemma}{lemmas}
\Crefname{lemma}{Lemma}{Lemmas}
\crefname{corollary}{corollary}{corollaries}
\Crefname{corollary}{Corollary}{Corollaries}
\crefname{remark}{remark}{remarks}
\Crefname{remark}{Remark}{Remarks}

% Bold notation from the supplied preamble.  Only symbols used below are kept.
\newcommand{\gb}{\mathbf{g}}
\newcommand{\ub}{\mathbf{u}}
\newcommand{\vb}{\mathbf{v}}
\newcommand{\wb}{\mathbf{w}}
\newcommand{\zb}{\mathbf{z}}
\newcommand{\Ab}{\mathbf{A}}
\newcommand{\Gb}{\mathbf{G}}
\newcommand{\Ib}{\mathbf{I}}
\newcommand{\Sb}{\mathbf{S}}
\newcommand{\Ub}{\mathbf{U}}
\newcommand{\Vb}{\mathbf{V}}
\newcommand{\Wb}{\mathbf{W}}
\newcommand{\Xb}{\mathbf{X}}
\newcommand{\Yb}{\mathbf{Y}}
\newcommand{\Zb}{\mathbf{Z}}

\newcommand{\EE}{\mathbb{E}}
\newcommand{\PP}{\mathbb{P}}
\newcommand{\RR}{\mathbb{R}}
\newcommand{\eps}{\varepsilon}
\newcommand{\norm}[1]{\left\lVert #1\right\rVert_2}
\newcommand{\normF}[1]{\left\lVert #1\right\rVert_{\mathrm F}}
\newcommand{\normop}[1]{\left\lVert #1\right\rVert_{\mathrm{op}}}
\newcommand{\HF}{H_{\mathrm F}}
\newcommand{\Hop}{H_{\mathrm{op}}}
\DeclareMathOperator{\polar}{polar}

\title{Provable Benefit of Spectral Optimizer over Stable Gradient Descent in Gaussian Associative Memory}
\date{}

\begin{document}
\maketitle

\begin{abstract}
We study optimization of an over-capacity Gaussian associativ -memory loss with Zipf weights, in the regime $N=\operatorname{poly}(d)$, $N\gtrsim d^2$, and $\alpha>1$. Over the accuracy range $\frac{(\log N)^\alpha}{d^{2(\alpha-1)}} \lesssim \varepsilon \lesssim \frac{\log N}{d^{\alpha-1}},$ we first show that the intrinsic Frobenius and operator-norm radii satisfy $R_{\mathrm F}(\varepsilon)^2\asymp d\,R_{\mathrm{op}}(\varepsilon)^2$. Since the corresponding global smoothness constants are both of constant order, this yields a factor-$d$ separation between the standard Frank-Wolfe smoothness-radius certificates in the two geometries. More importantly, we prove an actual optimization-time separation. For any deterministic, time-varying stable GD learning-rate schedule prescribed before the Gaussian sample, with $0\leq\eta_t\leq3$, we establish a uniform lower bound on its learning curve, while a target-free Muon-Frank-Wolfe trajectory satisfies a matching faster upper bound. Consequently, throughout the same accuracy range, $\tau_{\mathrm{GD}}(\varepsilon) \gtrsim d\,\tau_{\mathrm{Muon}}(\varepsilon).$ The GD lower bound also extends to weight decay under the corresponding stability condition. Finally, we show that the Muon learning-curve order is unchanged, up to constants, when its effective spectral learning rate is capped by an arbitrary prescribed constant $\bar\eta>0$. Thus the separation cannot be attributed merely to allowing Muon to take large effective steps.
\end{abstract}

% ============================================================
\section{Setup and main results}
% ============================================================

Fix $\alpha>1$.  Assume throughout that $N=\operatorname{poly}(d)$ and
$N\gtrsim d^2$.  All asymptotics are as $d\to\infty$.  Hidden constants
may depend on $\alpha$ and on the prescribed failure exponent, but not on
$d,N,\eps$, or the prescribed GD schedule.

Write $[N]:=\{1,\ldots,N\}$ and draw two mutually independent families
\[
 \ub_1,\ldots,\ub_N,
 \quad
 \vb_1,\ldots,\vb_N
 \overset{\mathrm{i.i.d.}}{\sim}
 \mathcal N\!\left(\mathbf 0,\frac1d\Ib_d\right).
\]
Let $p_i:=i^{-\alpha}/\sum_{k=1}^Nk^{-\alpha}$.  For
$\Wb\in\RR^{d\times d}$, define
\begin{align*}
 \ell_i(\Wb)
 &:={}
 \log\!\left(
  \sum_{j=1}^N
  \exp\!\left((\ub_j-\ub_i)^\top\Wb\vb_i\right)
 \right),\\
 L(\Wb)&:=\sum_{i=1}^Np_i\ell_i(\Wb).
\end{align*}
The term $j=i$ gives $L\geq0$, and $L(\mathbf 0)=\log N$.  Each $\ell_i$ is
a log-sum-exp of affine functions of $\Wb$, so each $\ell_i$ is convex.
Since the weights are positive, their weighted sum $L$ is also convex.

For $X\in\{\mathrm F,\mathrm{op}\}$, define the intrinsic radius
\[
 R_X(\eps)
 :=
 \inf\bigl\{\|\Wb\|_X:L(\Wb)\leq\eps\bigr\}.
\]
Whenever this set is nonempty, the infimum is attained: a minimizing
sequence is bounded, hence has a convergent subsequence, and the sublevel
set is closed because $L$ is continuous.  The global smoothness constants
are
\begin{align*}
 \HF
 &:={}
 \sup_{\Wb\in\RR^{d\times d}}
 \sup_{\Ab\ne\mathbf 0}
 \frac{\nabla^2L(\Wb)[\Ab,\Ab]}{\normF{\Ab}^2},\\
 \Hop
 &:={}
 \sup_{\Wb\in\RR^{d\times d}}
 \sup_{\Ab\ne\mathbf 0}
 \frac{\nabla^2L(\Wb)[\Ab,\Ab]}{\normop{\Ab}^2}.
\end{align*}

% \paragraph{Frank-wolfe.} For $X\in\{\mathrm F,\mathrm{op}\}$, Frank-Wolfe for target $\eps$ is
% run on the ball of radius $R_X(\eps/2)$.  Starting from
% $\Wb_0^X=\mathbf 0$, let $\gamma_t=2/(t+2)$ and define
% \[
%  \Sb_t^X
%  \in\mathop{\arg\min}_{\|\Sb\|_X\leq R_t}
%  \left\langle\nabla L(\Wb_t^X),\Sb\right\rangle_{\mathrm F},
%  \quad
%  \Wb_{t+1}^X=(1-\gamma_t)\Wb_t^X+\gamma_t\Sb_t^X.
% \]
% The operator-ball oracle is
% $-R_{\mathrm{op}}(\eps/2)\polar(\nabla L)$ and is called
% Muon-Frank-Wolfe.  The Frobenius-ball oracle is
% $-R_{\mathrm F}(\eps/2)\nabla L/\normF{\nabla L}$ and is called
% NGD-Frank-Wolfe; the fraction is zero when the gradient is zero.  Write
% \[
%  \tau_X(\eps):=\inf\{t\in\mathbb N:L(\Wb_t^X)\leq\eps\},
%  \quad
%  \tau_{\mathrm{Muon}}(\eps):=\tau_{\mathrm{op}}(\eps),
%  \quad
%  \tau_{\mathrm{NGD}}(\eps):=\tau_{\mathrm F}(\eps).
% \]
% If $\Gb=\Ub\Sigma\Vb^\top$ is a compact singular-value decomposition,
% define $\polar(\Gb):=\Ub\Vb^\top$, and set
% $\polar(\mathbf0):=\mathbf0$.  Then
% $\normop{\polar(\Gb)}\leq1$ and
% $\langle\Gb,\polar(\Gb)\rangle_{\mathrm F}=\|\Gb\|_*$.
\paragraph{Frank-Wolfe.}
For $X\in\{\mathrm F,\mathrm{op}\}$, let
$\{R_t\}_{t\geq0}$ be a prescribed nonnegative radius schedule.
Starting from $\Wb_0^X=\mathbf 0$, let $\{\gamma_t\}_{t\geq0}\subset(0,1]$
be a prescribed mixing schedule and define
\[
 \Sb_t^X\in \mathop{\arg\min}_{\|\Sb\|_X\leq R_t} \left\langle
 \nabla L(\Wb_t^X),\Sb
 \right\rangle_{\mathrm F},
 \quad
 \Wb_{t+1}^X=(1-\gamma_t)\Wb_t^X+\gamma_t\Sb_t^X.
\]
The corresponding linear minimization oracles are
$\Sb_t^{\mathrm{op}}=-R_t\polar\left(\nabla L(\Wb_t^{\mathrm{op}})\right)$
for the operator geometry, and
$\Sb_t^{\mathrm F}=-R_t\nabla L(\Wb_t^{\mathrm F})/\normF{\nabla L(\Wb_t^{\mathrm F})}
$
for the Frobenius geometry, with the latter fraction interpreted as zero
when the gradient is zero. For any target $0<\eps<\log N$, define hitting time:
\[
 \tau_X(\eps)
 :=
 \inf\left\{
 t\in\mathbb N:
 L(\Wb_t^X)\leq\eps
 \right\}.
\]
If $\Gb=\Ub\Sigma\Vb^\top$ is a compact singular-value decomposition,
define $\polar(\Gb):=\Ub\Vb^\top$, and set
$\polar(\mathbf0):=\mathbf0$.  Then
$\normop{\polar(\Gb)}\leq1$ and
$\langle\Gb,\polar(\Gb)\rangle_{\mathrm F}=\|\Gb\|_*$.

% \paragraph{Stable GD.}
% For a target $0<\varepsilon<\log N$, prescribe a deterministic sequence
% $\{\eta_t(\varepsilon)\}_{t\geq0}$ before drawing the Gaussian sample.  The
% sequence may depend arbitrarily on $(t,\varepsilon,d,N)$ and the fixed model
% parameters, and must only satisfy
% \begin{equation}
%  0\leq\eta_t(\varepsilon)\leq 3
%  \quad\text{for every }t\geq0.
%  \label{eq:eta-cap}
% \end{equation}
% Starting from $\bm W^{\mathrm{GD}}_0=\bm 0$, define
% \begin{align*}
%  \bm W^{\mathrm{GD}}_{t+1}
%  &=
%  \bm W^{\mathrm{GD}}_t
%  -\eta_t(\varepsilon)\nabla L(\bm W^{\mathrm{GD}}_t),
%  \\
%  \tau_{\mathrm{GD}}(\varepsilon)
%  &:=
%  \inf\{t\in\mathbb N:L(\bm W^{\mathrm{GD}}_t)\leq\varepsilon\}.
% \end{align*}
% The fixed schedule is suppressed from the hitting-time notation.  As usual,
% $\inf\varnothing=+\infty$.  The pre-sample requirement is needed
% only for the sign-conditioning step in the lower bound; it still permits an
% arbitrary target-dependent schedule.  Data-adaptive line searches are not
% included in this definition.
\paragraph{Stable GD.}
Prescribe a deterministic learning-rate sequence
$\{\eta_t\}_{t\geq 0}$ before drawing the Gaussian sample.  The sequence may
depend arbitrarily on $(t,d,N)$ and the fixed model parameters, but not on
the realized sample, and must satisfy
\begin{equation*}
    0\leq \eta_t\leq 3
    \qquad\text{for every }t\geq 0.
    \label{eq:eta-cap}
\end{equation*}
Starting from $\bm W^{\mathrm{GD}}_0=\bm 0$, define the single GD trajectory
\begin{equation*}
    \bm W^{\mathrm{GD}}_{t+1}
    =
    \bm W^{\mathrm{GD}}_t
    -\eta_t\nabla L(\bm W^{\mathrm{GD}}_t).
    \label{eq:gd-update}
\end{equation*}

The pre-sample requirement is used only in the sign-conditioning argument
for the lower bound; it still allows an arbitrary time-varying stable
learning-rate schedule.  In particular, the GD trajectory is independent of
the evaluation target $\varepsilon$.  Data-adaptive line searches are not
included.

\subsection*{Main results}

\paragraph{Overview.}


\begin{theorem}[Muon and NGD Frank-Wolfe bounds]
\label{thm:fw}
Fix $q>0$. For the Frank-Wolfe recursion above, for each $X\in\{\mathrm F,\mathrm{op}\}$ and target $\eps$, take
$R_t\equiv R_X(\eps/2)$ and $\gamma_t=2/(t+2).$ There exist $C_{\mathrm{low}},c_{\mathrm{up}}>0$ such that,
with probability at least $1-O(d^{-q})$, simultaneously for every
$
 C_{\mathrm{low}}
 \frac{(\log N)^\alpha}{d^{2(\alpha-1)}}
 \leq\eps\leq
 c_{\mathrm{up}}
 \frac{\log N}{d^{\alpha-1}},
$
one has
\begin{align*}
 \tau_{\mathrm{Muon}}(\eps)\lesssim
 \frac{\Hop R_{\mathrm{op}}(\eps/2)^2}{\eps}, \quad
 \tau_{\mathrm{NGD}}(\eps)
 \lesssim
 \frac{\HF R_{\mathrm F}(\eps/2)^2}{\eps}.
\end{align*}
In particular, the two standard Frank-Wolfe certificates satisfy
\[
 \frac{\HF R_{\mathrm F}(\eps/2)^2}{\eps}
 \eqsim
 d\,\frac{\Hop R_{\mathrm{op}}(\eps/2)^2}{\eps}.
\]
\end{theorem}

\begin{theorem}[Target-free Muon loss upper bound]
\label{thm:Muon-upper bound}
Fix $q>0$, and set the explicit deterministic constant $C_\alpha:=288(\alpha+2)$.  Starting from $\Wb_0^{\mathrm{Muon}}=\mathbf 0$, define, for every $t\geq0$,
\[
\gamma_t=\frac{2}{t+2},
\quad R_t=C_\alpha(\log N)\max\left\{1,\left(
\frac{t+1}{(\log N)d^{\alpha-1}}
\right)^{1/(2\alpha)}
\right\},
\]
and run the operator-norm Frank--Wolfe recursion
\[
\Wb_{t+1}^{\mathrm{Muon}}=
(1-\gamma_t)\Wb_t^{\mathrm{Muon}}
-\gamma_tR_t
\polar\!\left(\nabla L(\Wb_t^{\mathrm{Muon}})\right).
\]
There exist constants $C_{\mathrm{low}},c_{\mathrm{up}}>0$ such that,
with probability at least $1-O(d^{-q})$, simultaneously for every integer
$t$ satisfying
\[
C_{\mathrm{low}}(\log N)d^{\alpha-1}
\leq t
\leq
c_{\mathrm{up}}
\frac{d^{2\alpha-1}}{(\log N)^{\alpha-1}},
\]
one has
\[
L(\Wb_t^{\mathrm{Muon}})
\lesssim
(\log N)
\left(
\frac{dt}{\log N}
\right)^{-(\alpha-1)/\alpha}.
\]
\end{theorem}

\begin{theorem}[GD lower bound and actual separation]
\label{thm:gd}
Fix $q>0$.  Let $\{\eta_t\}_{t\geq0}$ be any deterministic GD schedule prescribed before the Gaussian sample and satisfying $0\leq \eta_t\leq 3.$ There exist $C_{\mathrm{low}},c_{\mathrm{up}}>0$, which may be chosen consistently with \Cref{thm:Muon-upper bound}, such that, with probability at least $1-O(d^{-q})$, simultaneously for every
\[
C_{\mathrm{low}}(\log N)d^\alpha
\leq t
\leq
c_{\mathrm{up}}
\frac{d^{2\alpha-1}}{(\log N)^{\alpha-1}},
\]
one has
\[
L(\Wb_t^{\mathrm{GD}})\gtrsim(\log N)\left(\frac{t}{\log N}\right)^{-(\alpha-1)/\alpha}\gtrsim
d^{(\alpha-1)/\alpha}
L(\Wb_t^{\mathrm{Muon}}).
\]

After adjusting $C_{\mathrm{low}}$ and $c_{\mathrm{up}}$ if necessary,
equivalently, for every
\[
C_{\mathrm{low}}
\frac{(\log N)^\alpha}{d^{2(\alpha-1)}}
\le \eps\le
c_{\mathrm{up}}
\frac{\log N}{d^{\alpha-1}},
\]
the corresponding hitting times satisfy
\[
\tau_{\mathrm{GD}}(\eps)
\gtrsim
(\log N)
\left(\frac{\log N}{\eps}\right)^{\alpha/(\alpha-1)}
\gtrsim
d\,\tau_{\mathrm{Muon}}(\eps).
\]
\end{theorem}

\begin{corollary}[GD with weight decay]
\label{cor:gd-weight-decay}
Let $\{\eta_t\}_{t\geq0}$ and $\{\lambda_t\}_{t\geq0}$ be arbitrary
nonnegative deterministic schedules prescribed before the Gaussian sample,
and consider
\[
\Wb^{\mathrm{GD+WD}}_{t+1}
=
(1-\eta_t\lambda_t)\Wb^{\mathrm{GD+WD}}_t
-\eta_t\nabla L(\Wb^{\mathrm{GD+WD}}_t),
\qquad
\Wb^{\mathrm{GD+WD}}_0=\mathbf 0.
\]
Suppose that, for every $t\geq0$,
$\eta_t\left(\lambda_t+\frac23\right)\leq 2.$
Then all the conclusions of \Cref{thm:gd} remain valid with
$\Wb_t^{\mathrm{GD}}$ replaced by $\Wb_t^{\mathrm{GD+WD}}$. 
% In particular,
% with probability at least $1-O(d^{-q})$, simultaneously for every $t$ in the iteration range of \Cref{thm:gd},
% one has
% \[
% L(\Wb_t^{\mathrm{GD+WD}})
% \gtrsim
% (\log N)
% \left(\frac{t}{\log N}\right)^{-(\alpha-1)/\alpha}
% \gtrsim
% d^{(\alpha-1)/\alpha}
% L(\Wb_t^{\mathrm{Muon}}).
% \]
% Equivalently, throughout the corresponding accuracy range,
% \[
% \tau_{\mathrm{GD+WD}}(\eps)
% \gtrsim
% (\log N)
% \left(\frac{\log N}{\eps}\right)^{\alpha/(\alpha-1)}
% \gtrsim
% d\,\tau_{\mathrm{Muon}}(\eps).
% \]
\end{corollary}

\begin{theorem}[Muon with an arbitrarily capped effective learning rate]
\label{thm:capped}
Fix $q>0$ and any constant $\bar\eta>0$.  Keep the same radius schedule
$R_t$ as in \Cref{thm:Muon-upper bound}, but replace the mixing schedule by
\[
\gamma_t=\frac{2}{t+2+a},
\quad
a=\max\left\{
0,\frac{2C_\alpha\log N}{\bar\eta}-2
\right\}.
\]
Then, for all sufficiently large $d$,
\[
\gamma_tR_t\leq\bar\eta,
\quad
\normop{\Wb_{t+1}^{\mathrm{Muon}}-\Wb_t^{\mathrm{Muon}}}
\leq2\bar\eta.
\]
Moreover, with
probability at least $1-O(d^{-q})$, simultaneously for every $t$ in the iteration range of \Cref{thm:Muon-upper bound} and \Cref{thm:gd},
one has
\[
L(\Wb_t^{\mathrm{Muon}})
\lesssim_{\bar\eta}
(\log N)\left(\frac{dt}{\log N}\right)^{-(\alpha-1)/\alpha}.
\]
\end{theorem}

% \begin{remark}[What the two separations mean]
% The factor $d$ in \Cref{thm:fw} compares the standard Frank-Wolfe
% smoothness-radius certificates; it is not a lower bound on the actual NGD
% hitting time.  The factor $d$ in \Cref{thm:gd} is an actual hitting-time
% separation.  In \Cref{thm:capped}, the quantity being capped is the
% effective learning rate $R_{\mathrm{op}}(\eps/2)\gamma_t$, not merely the mixing coefficient $\gamma_t\leq1$.
% \end{remark}

% ============================================================
\section{Proof of \Cref{thm:fw}}
% ============================================================

\subsection{Gaussian regularity and global smoothness}

\begin{lemma}[Gaussian regularity]
\label{lem:regularity-v2}
With probability at least $1-e^{-\Omega(d)}$, all the following hold:
\begin{enumerate}[label=\textup{(\roman*)},leftmargin=2.2em]
\item
\[
 \max_{i\leq N}
 \left\{
  \left|\norm{\ub_i}^2-1\right|,
  \left|\norm{\vb_i}^2-1\right|
 \right\}\leq\frac1{10},
\]
and
\[
 \max_{i\ne j}
 \left\{
  |\ub_i^\top\ub_j|,
  |\vb_i^\top\vb_j|
 \right\}\leq\frac1{100}.
\]
\item Every consecutive $B\subset[N]$ with
$d\leq|B|\leq2d$ and $\min B\geq d/3$ satisfies
\[
 \normop{[\vb_i]_{i\in B}}\leq4.
\]
\end{enumerate}
This event is unchanged by arbitrary sign flips of the vectors $\vb_i$.
\end{lemma}

\begin{proof}
For $\gb\sim\mathcal N(\mathbf0,\Ib_d/d)$, the chi-square tail bound gives
\[
 \PP\!\left(\left|\norm{\gb}^2-1\right|>\frac1{10}\right)
 \leq2e^{-c d}.
\]
For independent $\gb,\widetilde\gb$, conditionally on $\gb$ one has
$\gb^\top\widetilde\gb\sim
\mathcal N(0,\norm{\gb}^2/d)$.  Hence
\begin{align*}
 \PP\!\left(|\gb^\top\widetilde\gb|>\frac1{100}\right)
 &\leq
 \PP(\norm{\gb}>2)
 +\PP\!\left(
  |\gb^\top\widetilde\gb|>\frac1{100},\ \norm{\gb}\leq2
 \right)\\
 &\leq
 2e^{-c d}+2\exp\!\left(-\frac{d}{8\cdot10^4}\right)\\
 &\leq2e^{-c'd}.
\end{align*}
There are polynomially many vectors and pairs, so a union bound proves
part~(i).

If $|B|=m\in[d,2d]$, then
$[\vb_i]_{i\in B}=\Gb/\sqrt d$ for a standard Gaussian $d\times m$
matrix.  The Gaussian singular-value inequality gives
\begin{align*}
 \PP\!\left(\normop{\Gb/\sqrt d}>4\right)
 &\leq
 \PP\!\left(\normop{\Gb}>\sqrt d+\sqrt m+\sqrt d\right)\\
 &\leq2e^{-d/2},
\end{align*}
because $2+\sqrt{m/d}<4$.  There are polynomially many consecutive blocks,
so another union bound proves part~(ii).  Norms, absolute inner products,
and singular values are unchanged by column sign flips.
\end{proof}

\begin{lemma}[Sharp global smoothness]
\label{lem:smoothness-v2}
On the event in \Cref{lem:regularity-v2},
\[
 1\lesssim\HF\leq\Hop
 \leq\frac{1221}{2000}<\frac23.
\]
Consequently, every GD map
$\Wb\mapsto\Wb-\eta\nabla L(\Wb)$ with $0\leq\eta\leq3$ is
nonexpansive in Frobenius norm.
\end{lemma}

\begin{proof}
We first derive the Hessian without invoking the log-sum-exp variance formula
as a black box.  Define
\[
 \pi_{ji}(\Wb)
 :=
 \frac{e^{(\ub_j-\ub_i)^\top\Wb\vb_i}}
 {\sum_{k=1}^Ne^{(\ub_k-\ub_i)^\top\Wb\vb_i}}.
\]
Directional differentiation gives
\begin{align*}
 D\ell_i(\Wb)[\Ab]
 &=\sum_{j=1}^N\pi_{ji}(\Wb)
 (\ub_j-\ub_i)^\top\Ab\vb_i,\\
 D\pi_{ji}(\Wb)[\Ab]
 &=\pi_{ji}(\Wb)\Bigg[
 (\ub_j-\ub_i)^\top\Ab\vb_i-
 \sum_{k=1}^N\pi_{ki}(\Wb)
 (\ub_k-\ub_i)^\top\Ab\vb_i
 \Bigg].
\end{align*}
Substituting the second identity into the derivative of the first gives
\begin{align*}
 \nabla^2\ell_i(\Wb)[\Ab,\Ab]
 &=\sum_{j=1}^N\pi_{ji}(\Wb)
 \bigl((\ub_j-\ub_i)^\top\Ab\vb_i\bigr)^2-
 \left[
  \sum_{j=1}^N\pi_{ji}(\Wb)
  (\ub_j-\ub_i)^\top\Ab\vb_i
 \right]^2\\
 &=\operatorname{Var}_{j\sim\pi_{\cdot i}(\Wb)}
 \!\left((\ub_j-\ub_i)^\top\Ab\vb_i\right)\\
 &=\operatorname{Var}_{j\sim\pi_{\cdot i}(\Wb)}
 \!\left(\ub_j^\top\Ab\vb_i\right).
\end{align*}
The last equality holds because $\ub_i^\top\Ab\vb_i$ is constant in $j$.

If a scalar variable is supported on $[m,M]$, then
\begin{align*}
 \operatorname{Var}(Y)
 &=\min_{c\in\RR}\EE[(Y-c)^2]\\
 &\leq\EE\!\left[\left(Y-\frac{m+M}{2}\right)^2\right]
 \leq\frac{(M-m)^2}{4}.
\end{align*}
For $Y_j=\ub_j^\top\Ab\vb_i$, its range satisfies
\begin{align*}
 \max_jY_j-\min_jY_j
 &=\max_{j,k}(\ub_j-\ub_k)^\top\Ab\vb_i\\
 &\leq
 \max_{j,k}\norm{\ub_j-\ub_k}\,\norm{\Ab\vb_i},
\end{align*}
where the last line is Euclidean Cauchy--Schwarz.  On the regularity event,
the distance is zero for $j=k$, while for $j\ne k$,
\begin{align*}
 \norm{\ub_j-\ub_k}^2
 &=\norm{\ub_j}^2+\norm{\ub_k}^2-2\ub_j^\top\ub_k\\
 &\leq\frac{11}{10}+\frac{11}{10}+\frac2{100}
 =\frac{111}{50}.
\end{align*}
Consequently,
\begin{align*}
 \nabla^2L(\Wb)[\Ab,\Ab]
 &=\sum_{i=1}^Np_i\nabla^2\ell_i(\Wb)[\Ab,\Ab]\\
 &\leq\frac{111}{200}
 \sum_{i=1}^Np_i\norm{\Ab\vb_i}^2\\
 &\leq\frac{111}{200}\normop{\Ab}^2
 \sum_{i=1}^Np_i\norm{\vb_i}^2\\
 &\leq\frac{111}{200}\frac{11}{10}\normop{\Ab}^2
 =\frac{1221}{2000}\normop{\Ab}^2.
\end{align*}
The third line used
$\norm{\Ab\vb_i}\leq\normop{\Ab}\norm{\vb_i}$.  Replacing this by
$\norm{\Ab\vb_i}\leq\normF{\Ab}\norm{\vb_i}$ proves the Frobenius
upper bound.  Since $\normop{\Ab}\leq\normF{\Ab}$ for every $\Ab$, the
definitions also give $\HF\leq\Hop$.

We also verify explicitly that these Hessian bounds imply the smoothness
inequalities used below.  For $X\in\{\mathrm F,\mathrm{op}\}$ and arbitrary
$\Wb,\Ab$, the one-dimensional fundamental theorem of calculus gives
\begin{align*}
 L(\Wb+\Ab)-L(\Wb)
 &=\int_0^1
 \left\langle\nabla L(\Wb+s\Ab),\Ab\right\rangle_{\mathrm F}\,ds\\
 &=\left\langle\nabla L(\Wb),\Ab\right\rangle_{\mathrm F}
 +\int_0^1(1-s)\nabla^2L(\Wb+s\Ab)[\Ab,\Ab]\,ds\\
 &\leq\left\langle\nabla L(\Wb),\Ab\right\rangle_{\mathrm F}
 +\int_0^1(1-s)H_X\|\Ab\|_X^2\,ds\\
 &=\left\langle\nabla L(\Wb),\Ab\right\rangle_{\mathrm F}
 +\frac{H_X}{2}\|\Ab\|_X^2.
\end{align*}

We now prove that the lower bound is constant rather than merely positive.
Define, only in this paragraph,
\[
 \wb
 :=\ub_1+\ub_2
 -\frac{\norm{\ub_1}^2-\norm{\ub_2}^2}
 {\norm{\ub_1-\ub_2}^2}(\ub_1-\ub_2).
\]
By construction,
$(\ub_1-\ub_2)^\top\wb=0$, and hence
$\ub_1^\top\wb=\ub_2^\top\wb$.  The regularity bounds give
\begin{align*}
 \left|
  \frac{\norm{\ub_1}^2-\norm{\ub_2}^2}
  {\norm{\ub_1-\ub_2}^2}
 \right|
 &\leq\frac{1/5}{9/5-1/50}<\frac18,\\
 \ub_1^\top\wb
 &\geq
 \left(\frac9{10}-\frac1{100}\right)
 -\frac18\left(\frac{11}{10}+\frac1{100}\right)
 >\frac7{10},\\
 |\ub_j^\top\wb|
 &\leq\frac2{100}+\frac18\frac2{100}<\frac1{40}
 \qquad(j\geq3).
\end{align*}
Thus keys $1$ and $2$ are the unique tied maximizers of
$j\mapsto\ub_j^\top\wb$.

For $s>0$, let
\[
 \Wb_s:=s\,\frac{\wb\vb_1^\top}{\norm{\vb_1}^2},
 \qquad
 \Ab:=\frac{(\ub_1-\ub_2)\vb_1^\top}
 {\norm{\ub_1-\ub_2}\norm{\vb_1}}.
\]
The rank-one norm identities give
\begin{align*}
 \normF{\Ab}
 &=\frac{\norm{\ub_1-\ub_2}\norm{\vb_1}}
 {\norm{\ub_1-\ub_2}\norm{\vb_1}}=1,\\
 \normop{\Ab}&=\normF{\Ab}=1.
\end{align*}
Moreover, $\Wb_s\vb_1=s\wb$.  Therefore, in the first query's softmax,
\begin{align*}
 \pi_{j1}(\Wb_s)
 &=\frac{e^{s\ub_j^\top\wb}}
 {\sum_{k=1}^Ne^{s\ub_k^\top\wb}},
\end{align*}
where the common factor $e^{-s\ub_1^\top\wb}$ cancels.  The strict gap
proved above implies
\[
 \pi_{11}(\Wb_s),\pi_{21}(\Wb_s)\longrightarrow\frac12,
 \qquad
 \pi_{j1}(\Wb_s)\longrightarrow0\quad(j\geq3).
\]
For every $s>0$, the definition of $\HF$, the identity
$\normF{\Ab}=1$, and nonnegativity of every query variance give
\begin{align*}
 \HF
 &\geq\nabla^2L(\Wb_s)[\Ab,\Ab]\\
 &=\sum_{i=1}^Np_i
 \operatorname{Var}_{j\sim\pi_{\cdot i}(\Wb_s)}
 (\ub_j^\top\Ab\vb_i)\\
 &\geq p_1
 \operatorname{Var}_{j\sim\pi_{\cdot1}(\Wb_s)}
 (\ub_j^\top\Ab\vb_1).
\end{align*}
Taking $s\to\infty$ and using the already proved convergence of
$\pi_{\cdot1}(\Wb_s)$ yields
\begin{align*}
 \HF
 &\geq\frac{p_1}{4}
 \bigl((\ub_1-\ub_2)^\top\Ab\vb_1\bigr)^2\\
 &=\frac{p_1}{4}\norm{\ub_1-\ub_2}^2\norm{\vb_1}^2
 \gtrsim1.
\end{align*}
Here $p_1\geq1/\zeta(\alpha)$ and both squared norms are bounded below on
the regularity event.  This proves the lower smoothness bound.

The nonexpansiveness claim follows from cocoercivity, proved in full in
\Cref{lem:cocoercive-v2}: since $L$ is $2/3$-smooth, every
$0\leq\eta\leq2/(2/3)=3$ is allowed.
\end{proof}

\subsection{The standard Frank--Wolfe factor}

\begin{proposition}[Standard Frank--Wolfe smoothness--radius factor]
\label{prop:fw-factor-v2}
Fix $X\in\{\mathrm F,\mathrm{op}\}$ and $R>0$.  Frank--Wolfe on
$\{\Wb:\|\Wb\|_X\leq R\}$, started from zero with
$\gamma_t=2/(t+2)$, satisfies, for every $t\geq1$,
\[
 L(\Wb_t)-\min_{\|\Wb\|_X\leq R}L(\Wb)
 \leq\frac{8H_XR^2}{t+2}.
\]
Consequently, if $R\geq R_X(\eps/2)$, then its hitting time for $\eps$
satisfies
\[
 \inf\{t:L(\Wb_t)\leq\eps\}
 \lesssim1+\frac{H_XR^2}{\eps}.
\]
\end{proposition}

\begin{proof}
Let $\Wb_R^*$ minimize $L$ on the ball and write
$h_t:=L(\Wb_t)-L(\Wb_R^*)$.  Because $\gamma_t\in[0,1]$ and both
$\Wb_0$ and every oracle point $\Sb_t$ lie in the ball, induction gives
$\|\Wb_t\|_X\leq R$ for every $t$.  The exact linear oracle gives
\begin{align*}
 \left\langle\nabla L(\Wb_t),\Sb_t-\Wb_t\right\rangle_{\mathrm F}
 &=
 \min_{\|\Sb\|_X\leq R}
 \left\langle\nabla L(\Wb_t),\Sb\right\rangle_{\mathrm F}
 -\left\langle\nabla L(\Wb_t),\Wb_t\right\rangle_{\mathrm F}\\
 &\leq
 \left\langle\nabla L(\Wb_t),\Wb_R^*\right\rangle_{\mathrm F}
 -\left\langle\nabla L(\Wb_t),\Wb_t\right\rangle_{\mathrm F}\\
 &=\left\langle
 \nabla L(\Wb_t),\Wb_R^*-\Wb_t
 \right\rangle_{\mathrm F}.
\end{align*}
Convexity at $\Wb_t$ states
\[
 L(\Wb_R^*)
 \geq L(\Wb_t)
 +\left\langle\nabla L(\Wb_t),\Wb_R^*-\Wb_t\right\rangle_{\mathrm F}.
\]
Rearranging this inequality and substituting it into the preceding display
gives
\[
 \left\langle\nabla L(\Wb_t),\Sb_t-\Wb_t\right\rangle_{\mathrm F}
 \leq L(\Wb_R^*)-L(\Wb_t)=-h_t.
\]

The update gives the exact displacement
$\Wb_{t+1}-\Wb_t=\gamma_t(\Sb_t-\Wb_t)$.  The triangle inequality and
feasibility give
\[
 \|\Sb_t-\Wb_t\|_X
 \leq\|\Sb_t\|_X+\|\Wb_t\|_X
 \leq2R.
\]
Apply the smoothness inequality derived in \Cref{lem:smoothness-v2} with
$\Ab=\gamma_t(\Sb_t-\Wb_t)$ and then use the last two displays to obtain
\begin{align*}
 h_{t+1}
 &\leq h_t
 +\gamma_t\left\langle
  \nabla L(\Wb_t),\Sb_t-\Wb_t
 \right\rangle_{\mathrm F}
 +\frac{H_X\gamma_t^2}{2}\|\Sb_t-\Wb_t\|_X^2\\
 &\leq(1-\gamma_t)h_t+2H_XR^2\gamma_t^2.
\end{align*}
Since $\gamma_0=1$, this recurrence gives $h_1\leq2H_XR^2$.  Suppose
$h_t\leq8H_XR^2/(t+2)$.  Substituting $\gamma_t=2/(t+2)$ gives
\begin{align*}
 h_{t+1}
 &\leq
 \frac{t}{t+2}\frac{8H_XR^2}{t+2}
 +2H_XR^2\frac4{(t+2)^2}\\
 &=\frac{8H_XR^2(t+1)}{(t+2)^2}\\
 &\leq\frac{8H_XR^2}{t+3},
\end{align*}
where the last inequality is equivalent to
$(t+1)(t+3)\leq(t+2)^2$.  The base case follows from
$2H_XR^2\leq8H_XR^2/3$, so induction proves the first claim.

If $R\geq R_X(\eps/2)$, the attainment observation in the setup gives
$L(\Wb_R^*)\leq\eps/2$.  When
$8H_XR^2/(t+2)\leq\eps/2$, the first claim gives
\[
 L(\Wb_t)=L(\Wb_R^*)+h_t
 \leq\frac\eps2+\frac\eps2=\eps.
\]
Solving the preceding scalar inequality for $t$ proves the hitting bound.
\end{proof}

\subsection{The Gaussian comparator and prefix bounds}

\begin{proposition}[Uniform Gaussian prefix]
\label{prop:prefix-v2}
Fix $q>0$.  There is $c>0$ such that, with probability at least
$1-O(d^{-q})$, simultaneously for every integer
$d\leq K\leq c d^2/\log N$, the matrix
\[
 \Sb_K:=\sum_{k=1}^K\ub_k\vb_k^\top,
 \qquad
 \widehat\Wb_K:=16(\alpha+2)(\log N)\Sb_K
\]
satisfies
\begin{align*}
 \normop{\Sb_K}&\lesssim1+\sqrt{K/d},\\
 \normF{\Sb_K}&\lesssim\sqrt K,\\
 \left\|
  \sum_{i=1}^K
  \left(\ub_i-\frac1N\sum_{j=1}^N\ub_j\right)\vb_i^\top
 \right\|_{\mathrm F}
 &\lesssim\sqrt K,\\
 (\ub_i-\ub_j)^\top\Sb_K\vb_i
 &\geq\frac14
 \qquad(i\leq K,\ j\ne i),\\
 L(\widehat\Wb_K)&\lesssim(\log N)K^{1-\alpha}.
\end{align*}
\end{proposition}

\begin{proof}
We prove every line separately.

\paragraph{Operator and Frobenius norms.}
Let $\Ub_K:=[\ub_1,\ldots,\ub_K]$.  The Gaussian singular-value bound gives
\[
 \normop{\Ub_K}\leq2+\sqrt{K/d}
\]
except with probability $2e^{-d/2}$.  Conditional on $\Ub_K$, for every
standard basis vector $\mathbf e_r$,
\[
 \Sb_K\mathbf e_r
 =\sum_{k=1}^K(\vb_k)_r\ub_k.
\]
The coordinates $(\vb_k)_r$ are independent
$\mathcal N(0,1/d)$ variables.  Therefore the columns
$\Sb_K\mathbf e_r$ are conditionally independent centered Gaussians with
covariance
\[
 \EE[\Sb_K\mathbf e_r(\Sb_K\mathbf e_r)^\top\mid\Ub_K]
 =\frac1d\sum_{k=1}^K\ub_k\ub_k^\top
 =\frac1d\Ub_K\Ub_K^\top.
\]
Consequently, conditionally on $\Ub_K$,
\[
 \Sb_K\overset{\mathrm d}=
 (\Ub_K\Ub_K^\top)^{1/2}\frac{\Gb}{\sqrt d},
\]
where $\Gb$ is a standard $d\times d$ Gaussian matrix.  On the event
$\normop{\Gb}\leq3\sqrt d$,
\begin{align*}
 \normop{\Sb_K}
 &\leq
 \normop{(\Ub_K\Ub_K^\top)^{1/2}}
 \frac{\normop{\Gb}}{\sqrt d}\\
 &=\normop{\Ub_K}\frac{\normop{\Gb}}{\sqrt d}\\
 &\leq3(2+\sqrt{K/d})
 \lesssim1+\sqrt{K/d}.
\end{align*}
For the Frobenius norm, submultiplicativity in the form
$\|AB\|_{\mathrm F}\leq\|A\|_{\mathrm F}\|B\|_{\mathrm{op}}$ gives
\begin{align*}
 \normF{\Sb_K}
 &\leq
 \normF{(\Ub_K\Ub_K^\top)^{1/2}}
 \frac{\normop{\Gb}}{\sqrt d}\\
 &=\normF{\Ub_K}\frac{\normop{\Gb}}{\sqrt d}\\
 &=\left(\sum_{k=1}^K\norm{\ub_k}^2\right)^{1/2}
 \frac{\normop{\Gb}}{\sqrt d}\\
 &\lesssim\sqrt K.
\end{align*}
The Gaussian bound for $\Gb$, the bound for $\Ub_K$, and the norm event in
\Cref{lem:regularity-v2} each fail with probability $e^{-\Omega(d)}$.
Taking a union bound over the polynomially many values of $K$ preserves this
order.

Next, the triangle inequality and the rank-one Frobenius identity give
\begin{align*}
 &\left\|
  \sum_{i=1}^K
  \left(\ub_i-\frac1N\sum_{j=1}^N\ub_j\right)\vb_i^\top
 \right\|_{\mathrm F}\leq
 \left\|\sum_{i=1}^K\ub_i\vb_i^\top\right\|_{\mathrm F}
 +\left\|\frac1N\sum_{j=1}^N\ub_j\right\|
  \left\|\sum_{i=1}^K\vb_i\right\|.
\end{align*}
The two Gaussian sums satisfy
\begin{align*}
 \frac1N\sum_{j=1}^N\ub_j
 &\sim\mathcal N\!\left(\mathbf0,\frac1{dN}\Ib_d\right),\\
 \sum_{i=1}^K\vb_i
 &\sim\mathcal N\!\left(\mathbf0,\frac Kd\Ib_d\right).
\end{align*}
Chi-square concentration and a union bound over $K$ therefore give
\[
 \left\|\frac1N\sum_{j=1}^N\ub_j\right\|\lesssim\frac1{\sqrt N},
 \qquad
 \left\|\sum_{i=1}^K\vb_i\right\|\lesssim\sqrt K.
\]
Substitution into the triangle inequality, together with
$\normF{\Sb_K}\lesssim\sqrt K$, proves the centered-prefix bound.

\paragraph{Uniform margins.}
Fix $K$, $i\leq K$, and $j\ne i$.  Expanding every part of the sum gives
\begin{align*}
 (\ub_i-\ub_j)^\top\Sb_K\vb_i
 &=
 (\norm{\ub_i}^2-\ub_j^\top\ub_i)\norm{\vb_i}^2\\
 &\quad+
 \mathbf{1}_{\{j\leq K\}}
 (\ub_i^\top\ub_j-\norm{\ub_j}^2)(\vb_j^\top\vb_i)\\
 &\quad+
 \sum_{\substack{k\leq K\\k\notin\{i,j\}}}
 (\ub_i-\ub_j)^\top\ub_k(\vb_k^\top\vb_i).
\end{align*}
On the regularity event, the first line obeys
\[
 (\norm{\ub_i}^2-\ub_j^\top\ub_i)\norm{\vb_i}^2
 \geq\left(\frac9{10}-\frac1{100}\right)\frac9{10}
 >\frac45,
\]
while the second line has absolute value at most
\[
 \left(\frac{11}{10}+\frac1{100}\right)\frac1{100}<\frac3{200}.
\]

Condition on $(\ub_i,\ub_j,\vb_i)$, restricting first to conditioned values
that satisfy the three norm bounds in \Cref{lem:regularity-v2}.  This
restriction is measurable with respect to the conditioning.  For
$k\notin\{i,j\}$, the two factors in the last line are independent centered
Gaussians, because the first depends only on $\ub_k$ and the second only on
$\vb_k$.  For every such conditioned value, their conditional variances are
\begin{align*}
 \operatorname{Var}((\ub_i-\ub_j)^\top\ub_k
 \mid\ub_i,\ub_j)
 &=\frac{\norm{\ub_i-\ub_j}^2}{d}\lesssim\frac1d,\\
 \operatorname{Var}(\vb_k^\top\vb_i\mid\vb_i)
 &=\frac{\norm{\vb_i}^2}{d}\lesssim\frac1d.
\end{align*}
Thus their product is centered subexponential with conditional
subexponential norm $O(d^{-1})$.  Bernstein's inequality yields, for every
$s\geq1$,
\[
 \PP\!\left(
  \left|
   \sum_{\substack{k\leq K\\k\notin\{i,j\}}}
   (\ub_i-\ub_j)^\top\ub_k(\vb_k^\top\vb_i)
  \right|
  >C\left(\frac{\sqrt{Ks}}d+\frac{s}{d}\right)
  \,\middle|\,\ub_i,\ub_j,\vb_i
 \right)
 \leq2e^{-s}.
\]
For fixed $(K,i,j)$, let $\mathcal F_{Kij}$ denote the deviation event
inside this probability.  The tower property gives
\begin{align*}
 &\PP\!\left(
  \left\{
   \max\{\norm{\ub_i}^2,\norm{\ub_j}^2,\norm{\vb_i}^2\}
   \leq\frac{11}{10}
  \right\}
  \cap\mathcal F_{Kij}
 \right)\\
 &\qquad=
 \EE\!\left[
  \mathbf{1}_{\{
   \max\{\norm{\ub_i}^2,\norm{\ub_j}^2,\norm{\vb_i}^2\}\leq11/10
  \}}
  \PP(\mathcal F_{Kij}\mid\ub_i,\ub_j,\vb_i)
 \right]\\
 &\qquad\leq2e^{-s}.
\end{align*}
Take $s$ to be a sufficiently large fixed multiple of $\log N$.  A union
bound over all triples $(K,i,j)$, followed by adding the failure probability
of the norm event, gives total failure $O(d^{-q})$.  Since
$K\leq c d^2/\log N$, the displayed deviation is at most
$O(\sqrt c)+o(1)$.  Choose the fixed $c$ sufficiently small.  The last two
lines of the margin expansion then have total absolute value at most $1/2$,
and the margin is at least $4/5-1/2>1/4$.

\paragraph{Comparator loss.}
For $i\leq K$, the margin gives
\begin{align*}
 \ell_i(\widehat\Wb_K)
 &=\log\!\left[
  1+\sum_{j\ne i}
  \exp\!\left(
   -16(\alpha+2)(\log N)
   (\ub_i-\ub_j)^\top\Sb_K\vb_i
  \right)
 \right]\\
 &\leq\log\!\left(1+(N-1)N^{-4(\alpha+2)}\right)\\
 &\leq N^{-4\alpha-7}.
\end{align*}

For $i>K$, condition on all key vectors and on $\vb_1,\ldots,\vb_K$.
Then $\widehat\Wb_K$ is fixed and $\vb_i$ remains Gaussian.  Define only
inside this paragraph
\[
 f_i(\vb)
 :=\log\sum_{j=1}^N
 \exp\!\left((\ub_j-\ub_i)^\top\widehat\Wb_K\vb\right).
\]
Its gradient is
\[
 \nabla f_i(\vb)
 =\widehat\Wb_K^\top
 \left(\sum_{j=1}^N\pi_j(\vb)\ub_j-\ub_i\right),
\]
where the $\pi_j(\vb)$ are nonnegative and sum to one.  Hence
\begin{align*}
 \norm{\nabla f_i(\vb)}
 &\leq\normop{\widehat\Wb_K}
 \left(\sum_j\pi_j(\vb)\norm{\ub_j}+\norm{\ub_i}\right)\\
 &\leq3\normop{\widehat\Wb_K}.
\end{align*}
The product-space gradient of the weighted tail therefore satisfies
\begin{align*}
 \left\|
  \nabla_{(\vb_i)_{i>K}}
  \sum_{i>K}p_if_i(\vb_i)
 \right\|_2^2=
 \sum_{i>K}p_i^2\norm{\nabla f_i(\vb_i)}^2\leq
 9\normop{\widehat\Wb_K}^2\sum_{i>K}p_i^2.
\end{align*}
For every conditioning satisfying
$\max_j\|\ub_j\|^2\le 11/10$, the following holds.
The random vectors $(\vb_i)_{i>K}$ are independent and
$\vb_i\sim\mathcal N(\mathbf0,\Ib_d/d)$, so the product measure
is a centered Gaussian on the product space with covariance operator
$\sigma^2\Ib$ where $\sigma^2=1/d$.

Define the function
\[
g\bigl((\vb_i)_{i>K}\bigr)
:=\sum_{i>K}p_i\,f_i(\vb_i).
\]
Its Euclidean gradient (with respect to the product Euclidean structure)
satisfies
\begin{align*}
\|\nabla g\|_2^2
&=\sum_{i>K}p_i^2\|\nabla f_i(\vb_i)\|^2
\le\sum_{i>K}p_i^2\cdot\bigl(3\|\widehat\Wb_K\|_{\mathrm{op}}\bigr)^2
=9\|\widehat\Wb_K\|_{\mathrm{op}}^2\sum_{i>K}p_i^2,
\end{align*}
where the bound $\|\nabla f_i\|\le 3\|\widehat\Wb_K\|_{\mathrm{op}}$
holds uniformly on the conditioning event
$\max_j\|\ub_j\|^2\le 11/10$.
Consequently $g$ is Lipschitz continuous with constant
\[
L=3\|\widehat\Wb_K\|_{\mathrm{op}}\sqrt{\sum_{i>K}p_i^2}.
\]

The standard Gaussian concentration inequality for an $L$-Lipschitz
function of a Gaussian vector with covariance $\sigma^2\Ib$ asserts that
\[
\PP\bigl(|g-\EE[g\mid\text{conditioning}]|>t\bigm|\text{conditioning}\bigr)
\le 2\exp\Bigl(-\frac{t^2}{2L^2\sigma^2}\Bigr).
\]
Choosing
\[
t=C\|\widehat\Wb_K\|_{\mathrm{op}}
\sqrt{\frac{\bigl(\sum_{i>K}p_i^2\bigr)\log N}{d}}
\]
with a sufficiently large absolute constant $C$ (e.g.\ $C=3\sqrt{2(q+4)}$)
makes the right-hand side at most $N^{-q-4}$. 
Therefore, with conditional probability at least $1-N^{-q-4}$,
\begin{align*}
\biggl|
\sum_{i>K}p_i
\bigl(f_i(\vb_i)-\EE[f_i(\vb_i)\mid\text{conditioning}]\bigr)
\biggr|
&\lesssim
\|\widehat\Wb_K\|_{\mathrm{op}}
\sqrt{\frac{\bigl(\sum_{i>K}p_i^2\bigr)\log N}{d}}.
\end{align*}
For fixed $K$, let $\mathcal F_K$ be failure of this concentration
inequality.  Since the key-norm event is measurable with respect to the
conditioning, the tower property gives
\begin{align*}
 &\PP\!\left(
  \left\{\max_j\norm{\ub_j}^2\leq\frac{11}{10}\right\}
  \cap\mathcal F_K
 \right)\\
 &\qquad=
 \EE\!\left[
  \mathbf{1}_{\{\max_j\norm{\ub_j}^2\leq11/10\}}
  \PP(\mathcal F_K\mid\text{conditioning})
 \right]\\
 &\qquad\leq N^{-q-4}.
\end{align*}
There are at most $d^2$ admissible values of $K$.  A union bound, using
$N\gtrsim d^2$, and then adding the failure probability of the key-norm
event make this estimate simultaneous in $K$ with total failure
$O(d^{-q})$.

For the conditional mean, Jensen's inequality and the Gaussian
moment-generating function give every intermediate step:
\begin{align*}
 \EE[f_i(\vb_i)\mid\text{conditioning}]
 &\leq
 \log\EE\!\left[
  \sum_{j=1}^N
  e^{(\ub_j-\ub_i)^\top\widehat\Wb_K\vb_i}
  \,\middle|\,\text{conditioning}
 \right]\\
 &=\log\sum_{j=1}^N
 \EE\!\left[
  e^{(\ub_j-\ub_i)^\top\widehat\Wb_K\vb_i}
  \,\middle|\,\text{conditioning}
 \right]\\
 &=\log\sum_{j=1}^N
 \exp\!\left(
  \frac{\norm{\widehat\Wb_K^\top(\ub_j-\ub_i)}^2}{2d}
 \right)\\
 &\leq\log N
 +\max_j\frac{\norm{\widehat\Wb_K^\top(\ub_j-\ub_i)}^2}{2d}\\
 &\leq\log N+\frac3d\normop{\widehat\Wb_K}^2.
\end{align*}
The third line uses that the exponent is a centered Gaussian with variance
$\norm{\widehat\Wb_K^\top(\ub_j-\ub_i)}^2/d$.  The last line uses
$\norm{\ub_j-\ub_i}^2\leq111/50<6$.

Adding the head loss, the conditional tail mean, and the concentration term
gives
\begin{align*}
 L(\widehat\Wb_K)
 &\leq
 (\log N)\sum_{i>K}p_i
 +\frac3d\left(\sum_{i>K}p_i\right)
 \normop{\widehat\Wb_K}^2\\
 &\quad+
 O\!\left(
  \normop{\widehat\Wb_K}
  \sqrt{\frac{(\sum_{i>K}p_i^2)\log N}{d}}
 \right)
 +N^{-4\alpha-7}.
\end{align*}
Integral comparison gives
\[
 \sum_{i>K}p_i\lesssim Kp_K,
 \qquad
 \sum_{i>K}p_i^2\lesssim Kp_K^2.
\]
Using
$\normop{\widehat\Wb_K}\lesssim(\log N)(1+\sqrt{K/d})$, the conditional
mean error divided by $Kp_K\log N$ satisfies
\begin{align*}
 \frac{
  d^{-1}(\sum_{i>K}p_i)\normop{\widehat\Wb_K}^2
 }{Kp_K\log N}&\lesssim
 \frac{d^{-1}(Kp_K)(\log N)^2(1+K/d)}
 {Kp_K\log N}\\
 &=
 \frac{\log N}{d}+\frac{K\log N}{d^2}.
\end{align*}
The concentration error divided by the same quantity satisfies
\begin{align*}
 \frac{
  \normop{\widehat\Wb_K}
  \sqrt{(\sum_{i>K}p_i^2)\log N/d}
 }{Kp_K\log N}&\lesssim
 \frac{
  (\log N)(1+\sqrt{K/d})p_K\sqrt{K\log N/d}
 }{Kp_K\log N}\\
 &=\sqrt{\frac{\log N}{Kd}}+\frac{\sqrt{\log N}}d.
\end{align*}
Finally, since $Kp_K\eqsim K^{1-\alpha}$ and $K\leq d^2$,
\begin{align*}
 \frac{N^{-4\alpha-7}}{Kp_K\log N}
 &\lesssim
 \frac{K^{\alpha-1}}{N^{4\alpha+7}\log N}\\
 &\leq
 \frac{d^{2(\alpha-1)}}{N^{4\alpha+7}\log N}=o(1).
\end{align*}
The first ratio is $O(c)+o(1)$ and the second is $o(1)$ uniformly in the
admissible range, and the last display treats the head error.  Shrink the
fixed $c$ and then take $d$ large.  The loss estimate then yields
\begin{align*}
 L(\widehat\Wb_K)
 &\lesssim(\log N)\sum_{i>K}p_i+Kp_K\log N\\
 &\lesssim(\log N)K^{1-\alpha}.
\end{align*}
Intersecting all events and taking the stated union bounds proves the
proposition.
\end{proof}

\subsection{Intrinsic radius bounds and completion of the proof}

\begin{proof}[Proof of \Cref{thm:fw}]
Fix a target in the stated range.  Since
$1\leq\sum_{j=1}^Nj^{-\alpha}\leq\zeta(\alpha)$,
\begin{align*}
 \frac{K^{-\alpha}}{\zeta(\alpha)}
 &\leq p_K\leq K^{-\alpha},\\
 \sum_{i>K}p_i
 &=\frac{\sum_{i=K+1}^Ni^{-\alpha}}
 {\sum_{j=1}^Nj^{-\alpha}}\leq\sum_{i=K+1}^{\infty}i^{-\alpha}\leq\int_K^\infty y^{-\alpha}\,dy
 =\frac{K^{1-\alpha}}{\alpha-1}.
\end{align*}
If $K\leq N/2$, the reverse tail bound follows from the first $K$ terms:
\begin{align*}
 \sum_{i>K}p_i
 &\geq\sum_{i=K+1}^{2K}p_i\geq Kp_{2K}\geq\frac{2^{-\alpha}}{\zeta(\alpha)}K^{1-\alpha}.
\end{align*}
Consequently, uniformly for $K\leq N/2$,
$p_K\eqsim K^{-\alpha}$,
$\sum_{i>K}p_i\eqsim K^{1-\alpha}$, and
$Kp_K\eqsim K^{1-\alpha}$.

The target range gives the auxiliary scale directly, without introducing a
new symbol:
\begin{align*}
 \eps
 &\leq c_{\mathrm{up}}\frac{\log N}{d^{\alpha-1}}
 \quad\Longrightarrow\quad
 \left(\frac{\log N}{\eps}\right)^{1/(\alpha-1)}
 \gtrsim d,\\
 \eps
 &\geq C_{\mathrm{low}}
 \frac{(\log N)^\alpha}{d^{2(\alpha-1)}}
 \quad\Longrightarrow\quad
 \left(\frac{\log N}{\eps}\right)^{1/(\alpha-1)}
 \lesssim\frac{d^2}{\log N}.
\end{align*}
Moreover, $N\gtrsim d^2$ implies
$c d^2/\log N\leq N/2$ for all sufficiently large $d$.  Thus every
auxiliary $K$ used below lies in the range where the preceding Zipf
estimates are valid.

For the operator-radius upper bound, choose an integer
\[
 K\eqsim\left(\frac{\log N}{\eps}\right)^{1/(\alpha-1)}
\]
with a sufficiently large fixed implicit constant.  The target range,
after increasing $C_{\mathrm{low}}$ and decreasing $c_{\mathrm{up}}$ to
leave fixed-factor slack, therefore ensures
$d\leq K\leq c d^2/\log N$.  The last bound in
\Cref{prop:prefix-v2} then gives
\begin{align*}
 L(\widehat\Wb_K)
 &\lesssim(\log N)K^{1-\alpha}
 \leq\eps.
\end{align*}
Thus $\widehat\Wb_K$ is feasible in the definition of
$R_{\mathrm{op}}(\eps)$.  Since $K\geq d$,
\begin{align*}
 R_{\mathrm{op}}(\eps)
 &\leq\normop{\widehat\Wb_K}\\
 &\lesssim(\log N)(1+\sqrt{K/d})\\
 &\lesssim\frac{\log N}{\sqrt d}
 \left(\frac{\log N}{\eps}\right)^{1/[2(\alpha-1)]}.
\end{align*}

We next prove the nontrivial Frobenius-radius lower bound.  Choose an integer
of the same order, now with a sufficiently small fixed implicit constant,
so that
\[
 \eps\leq\frac18Kp_K\log N.
\]
This is possible because $Kp_K\log N\eqsim(\log N)K^{1-\alpha}$; the target
range again makes this $K$ admissible.  Let $\Wb$ be any matrix satisfying
$L(\Wb)\leq\eps$.  Since $p_i\geq p_K$ for $i\leq K$ and every loss is
nonnegative,
\begin{align*}
 p_K\sum_{i=1}^K\ell_i(\Wb)
 &\leq\sum_{i=1}^Kp_i\ell_i(\Wb)\leq L(\Wb)\leq\eps,
\end{align*}
and therefore
\[
 \sum_{i=1}^K\ell_i(\Wb)
 \leq\frac{\eps}{p_K}
 \leq\frac18K\log N.
\]

For every $i$, Jensen's inequality gives
\begin{align*}
 \ell_i(\Wb)
 &=\log N
 +\log\!\left[
  \frac1N\sum_{j=1}^N
  e^{(\ub_j-\ub_i)^\top\Wb\vb_i}
 \right]\\
 &\geq\log N
 +\frac1N\sum_{j=1}^N
 (\ub_j-\ub_i)^\top\Wb\vb_i\\
 &=\log N
 -\left(\ub_i-\frac1N\sum_{j=1}^N\ub_j\right)^\top\Wb\vb_i.
\end{align*}
Rearranging this inequality, summing from $1$ to $K$, and then using the
preceding loss bound gives
\begin{align*}
 \left\langle
  \Wb,
  \sum_{i=1}^K
  \left(\ub_i-\frac1N\sum_{j=1}^N\ub_j\right)\vb_i^\top
 \right\rangle_{\mathrm F} &=
 \sum_{i=1}^K
 \left(\ub_i-\frac1N\sum_{j=1}^N\ub_j\right)^\top\Wb\vb_i\\
 &\geq
 K\log N-\sum_{i=1}^K\ell_i(\Wb)\geq\frac78K\log N.
\end{align*}
Frobenius Cauchy--Schwarz and the centered-prefix bound in
\Cref{prop:prefix-v2} now give
\begin{align*}
 \frac78K\log N
 &\leq
 \left|
  \left\langle
   \Wb,
   \sum_{i=1}^K
   \left(\ub_i-\frac1N\sum_{j=1}^N\ub_j\right)\vb_i^\top
  \right\rangle_{\mathrm F}
 \right|\\
 &\leq
 \normF{\Wb}
 \left\|
  \sum_{i=1}^K
  \left(\ub_i-\frac1N\sum_{j=1}^N\ub_j\right)\vb_i^\top
 \right\|_{\mathrm F}\\
 &\lesssim\normF{\Wb}\sqrt K.
\end{align*}
Thus every feasible $\Wb$ satisfies
\[
 \normF{\Wb}
 \gtrsim(\log N)\sqrt K
 \eqsim(\log N)
 \left(\frac{\log N}{\eps}\right)^{1/[2(\alpha-1)]}.
\]
Taking the infimum proves the Frobenius lower bound.

For every $d\times d$ matrix,
$\normF{\Wb}\leq\sqrt d\normop{\Wb}$.  Taking infima over the same
sublevel set gives, with every step shown,
\begin{align*}
 R_{\mathrm F}(\eps)
 &=\inf_{L(\Wb)\leq\eps}\normF{\Wb}\\
 &\leq\inf_{L(\Wb)\leq\eps}\sqrt d\normop{\Wb}\\
 &=\sqrt d\,R_{\mathrm{op}}(\eps).
\end{align*}
Combining the Frobenius lower bound with
$R_{\mathrm F}(\eps)\leq\sqrt d\,R_{\mathrm{op}}(\eps)$ gives
\begin{align*}
 R_{\mathrm{op}}(\eps)
 &\geq\frac{R_{\mathrm F}(\eps)}{\sqrt d}\gtrsim
 \frac{\log N}{\sqrt d}
 \left(\frac{\log N}{\eps}\right)^{1/[2(\alpha-1)]}.
\end{align*}
Together with the operator-radius upper bound, this proves
\[
 R_{\mathrm{op}}(\eps)
 \eqsim
 \frac{\log N}{\sqrt d}
 \left(\frac{\log N}{\eps}\right)^{1/[2(\alpha-1)]}.
\]
Conversely, the operator-radius upper bound and the same deterministic norm
comparison give
\begin{align*}
 R_{\mathrm F}(\eps)
 &\leq\sqrt d\,R_{\mathrm{op}}(\eps)\lesssim
 (\log N)
 \left(\frac{\log N}{\eps}\right)^{1/[2(\alpha-1)]}.
\end{align*}
Together with the Frobenius-radius lower bound, this proves
\[
 R_{\mathrm F}(\eps)
 \eqsim
 (\log N)
 \left(\frac{\log N}{\eps}\right)^{1/[2(\alpha-1)]}.
\]
In particular,
\[
 R_{\mathrm F}(\eps)^2\eqsim dR_{\mathrm{op}}(\eps)^2.
\]
Choose $C_{\mathrm{low}}$ sufficiently large and
$c_{\mathrm{up}}$ sufficiently small so that the two preceding choices of
$K$ remain in $[d,c d^2/\log N]$ both for $\eps$ and for $\eps/2$.
Repeating the same two radius arguments with $\eps/2$ in place of $\eps$
therefore gives
\[
 R_X(\eps/2)^2\eqsim R_X(\eps)^2,
 \qquad X\in\{\mathrm F,\mathrm{op}\}.
\]

By operator-nuclear norm duality and Frobenius self-duality,
the linear minimization oracles specified in the setup are exact:
\[
 \min_{\normop{\Sb}\le R}\langle \Gb,\Sb\rangle_{\mathrm F}
 =-R\|\Gb\|_*,
 \qquad
 \min_{\normF{\Sb}\le R}\langle \Gb,\Sb\rangle_{\mathrm F}
 =-R\normF{\Gb}.
\]
The minima are attained by
$-R\polar(\Gb)$ and
$-R\Gb/\normF{\Gb}$, respectively, with the latter interpreted as
$\mathbf 0$ when $\Gb=\mathbf 0$.
Hence \Cref{prop:fw-factor-v2} applies with $R=R_X(\eps/2)$.


The hitting bound in \Cref{prop:fw-factor-v2}, the radius scales, and
$\HF,\Hop\eqsim1$ give
\begin{align*}
 \frac{\Hop R_{\mathrm{op}}(\eps/2)^2}{\eps}
 &\eqsim
 \frac{(\log N)^2}{d\eps}
 \left(\frac{\log N}{\eps}\right)^{1/(\alpha-1)}=\frac{\log N}{d}
 \left(\frac{\log N}{\eps}\right)^{\alpha/(\alpha-1)},\\
 \frac{\HF R_{\mathrm F}(\eps/2)^2}{\eps}
 &\eqsim
 \frac{(\log N)^2}{\eps}
 \left(\frac{\log N}{\eps}\right)^{1/(\alpha-1)}=(\log N)
 \left(\frac{\log N}{\eps}\right)^{\alpha/(\alpha-1)}.
\end{align*}
Both quantities diverge in the target range, so they absorb the additive
$1$.  Finally, $1\lesssim\HF\leq\Hop\lesssim1$ and the sharp radius relation
give
\begin{align*}
 \frac{\HF R_{\mathrm F}(\eps/2)^2}
 {\Hop R_{\mathrm{op}}(\eps/2)^2}
 &=\frac{\HF}{\Hop}
 \frac{R_{\mathrm F}(\eps/2)^2}{R_{\mathrm{op}}(\eps/2)^2}\\
 &\eqsim 1\cdot d=d.
\end{align*}
This is a comparison of the standard Frank-Wolfe smoothness-radius
certificates, not a lower bound on $\tau_{\mathrm{NGD}}(\eps)$.  This
completes the proof.
\end{proof}

% ============================================================

\section{Proof of \Cref{thm:Muon-upper bound}}
% ============================================================

We first extract from \Cref{prop:prefix-v2} a comparator at the scale
required by each iteration.  We then prove a deterministic
changing-radius Frank--Wolfe inequality and the two scalar summation bounds
needed to solve it.
Through this section, we write $M:=(\log N)d^{\alpha-1}.$
Since $N\gtrsim d^2$ and $\alpha>1$, one has $M\geq1$ for all
sufficiently large $d$.

\begin{lemma}[Time-dependent Gaussian comparators]
\label{lem:Muon-time-comparators}
Fix $q>0$, and let $C_\alpha=288(\alpha+2)$.  For every integer $s\geq0$,
define
\[
 K_s
 :=
 \left\lceil
  d\max\left\{
    1,
    \left(
    \frac{s+1}{M}
   \right)^{1/\alpha}
  \right\}
 \right\rceil.
\]
There exist constants $A_{\alpha,q},c_*>0$ such that, on the event in
\Cref{prop:prefix-v2}, simultaneously for every integer
\[
 1\leq t
 \leq
 c_*
 \frac{d^{2\alpha-1}}{(\log N)^{\alpha-1}}
\]
and every $0\leq s<t$, the integer $K_s$ lies in the admissible range of
\Cref{prop:prefix-v2} and satisfies
\begin{align*}
 \normop{\widehat\Wb_{K_s}}
 &\leq
 C_\alpha(\log N)
 \max\left\{
  1,
  \left(
   \frac{s+1}{M}
  \right)^{1/(2\alpha)}
 \right\},\\
 L(\widehat\Wb_{K_s})
 &\leq
 A_{\alpha,q}(\log N)d^{1-\alpha}
 \max\left\{
  1,
  \frac{s+1}{M}
 \right\}^{-(\alpha-1)/\alpha}.
\end{align*}
\end{lemma}
\begin{proof}
By \Cref{prop:prefix-v2}, on an event of probability at least
$1-O(d^{-q})$, there exist $A_{\alpha,q},c_0>0$ such that,
simultaneously for every integer
$d\leq K\leq c_0\frac{d^2}{\log N},$
The loss satisfies
\[
L(\widehat\Wb_K)
\leq
A_{\alpha,q}(\log N)K^{1-\alpha},
\]
and the explicit operator-norm estimate in the proof of \Cref{prop:prefix-v2} gives
\[
\normop{\widehat\Wb_K}
\leq
48(\alpha+2)(\log N)
\left(2+\sqrt{K/d}\right)
\leq
144(\alpha+2)(\log N)\sqrt{K/d},
\]
where the last inequality uses $K\geq d$.

From the definition of $K_s$,
\[
d\max\left\{
1,\left(\frac{s+1}{M}\right)^{1/\alpha}
\right\}
\leq K_s
\leq
2d\max\left\{
1,\left(\frac{s+1}{M}\right)^{1/\alpha}
\right\}.
\]
In particular, $K_s\geq d$.  Moreover, if
$t\leq c_*\frac{d^{2\alpha-1}}{(\log N)^{\alpha-1}}$ and $0\le s<t$,
then, using $M=(\log N)d^{\alpha-1}$,
\[
\left(\frac{s+1}{M}\right)^{1/\alpha}
\leq
c_*^{1/\alpha}\frac{d}{\log N}.
\]
Hence
\[
K_s
\leq
2d+2c_*^{1/\alpha}\frac{d^2}{\log N}.
\]
Choosing $c_*>0$ sufficiently small, and then $d$ sufficiently large,
ensures
\[
K_s\leq c_0\frac{d^2}{\log N},
\]
so every $K_s$ lies in the simultaneous range of
\Cref{prop:prefix-v2}.

Therefore, since $C_\alpha=288(\alpha+2)$, we have
\begin{align*}
\normop{\widehat\Wb_{K_s}}
&\leq
144(\alpha+2)(\log N)\sqrt{K_s/d}\\
&\leq
144\sqrt2(\alpha+2)(\log N)
\max\left\{
1,
\left(\frac{s+1}{M}\right)^{1/(2\alpha)}
\right\}\\
&\leq
C_\alpha(\log N)
\max\left\{
1,
\left(\frac{s+1}{M}\right)^{1/(2\alpha)}
\right\},
\end{align*}

Finally, because $1-\alpha<0$,
\begin{align*}
L(\widehat\Wb_{K_s})
&\leq
A_{\alpha,q}(\log N)K_s^{1-\alpha}\\
&\leq
A_{\alpha,q}(\log N)d^{1-\alpha}
\max\left\{
1,\frac{s+1}{M}
\right\}^{-(\alpha-1)/\alpha}.
\end{align*}
This proves the lemma.
\end{proof}

\begin{lemma}[Changing-radius Frank-Wolfe inequality]
\label{lem:Muon-changing-radius}
Let $R_0\leq R_1\leq\cdots$ and let
$\gamma_s=2/(s+2)$.  Starting from $\Wb_0=\mathbf0$, let
\begin{align*}
 \Sb_s\in\mathop{\arg\min}_{\normop{\Sb}\leq R_s}
 \left\langle\nabla L(\Wb_s),\Sb\right\rangle_{\mathrm F}, \quad\Wb_{s+1}
 &=(1-\gamma_s)\Wb_s+\gamma_s\Sb_s.
\end{align*}
Fix an integer $t\geq1$.  Suppose that, for every $0\leq s<t$, there is a
prefix comparator $\widehat\Wb_{K_s}$ satisfying
$\normop{\widehat\Wb_{K_s}}\leq R_s$.  Then
\begin{align*}
 t(t+1)L(\Wb_t)
 &\leq
 2\sum_{s=0}^{t-1}(s+1)L(\widehat\Wb_{K_s})
 +8\Hop
 \sum_{s=0}^{t-1}\frac{s+1}{s+2}R_s^2.
\end{align*}
\end{lemma}

\begin{proof}
We first show that the iterate at time $s$ lies in the current ball.  The
claim holds at $s=0$ because $\Wb_0=\mathbf0$.  If
$\normop{\Wb_s}\leq R_s$, then the oracle satisfies
$\normop{\Sb_s}\leq R_s$, and hence
\begin{align*}
 \normop{\Wb_{s+1}}
 &=\normop{(1-\gamma_s)\Wb_s+\gamma_s\Sb_s}\\
 &\leq
 (1-\gamma_s)\normop{\Wb_s}
 +\gamma_s\normop{\Sb_s}\\
 &\leq
 (1-\gamma_s)R_s+\gamma_sR_s\\
 &=R_s\leq R_{s+1}.
\end{align*}
Induction therefore gives, for every $s\geq0$,
\begin{align*}
 \normop{\Wb_s}&\leq R_s,\\
 \normop{\Sb_s-\Wb_s}
 &\leq
 \normop{\Sb_s}+\normop{\Wb_s}
 \leq2R_s.
\end{align*}

The exact linear minimization oracle and the feasibility of
$\widehat\Wb_{K_s}$ give
\begin{align}
\nonumber
 \left\langle
  \nabla L(\Wb_s),\Sb_s-\Wb_s
 \right\rangle_{\mathrm F}
 &=
 \min_{\normop{\Sb}\leq R_s}
 \left\langle\nabla L(\Wb_s),\Sb\right\rangle_{\mathrm F}
 -\left\langle\nabla L(\Wb_s),\Wb_s\right\rangle_{\mathrm F}\\
 &\leq
 \left\langle
  \nabla L(\Wb_s),\widehat\Wb_{K_s}-\Wb_s
 \right\rangle_{\mathrm F}.\label{eq:fw-comparater}
\end{align}
Convexity at $\Wb_s$ gives
\begin{align}
 L(\widehat\Wb_{K_s})\geq
 L(\Wb_s)
 +\left\langle
  \nabla L(\Wb_s),\widehat\Wb_{K_s}-\Wb_s
 \right\rangle_{\mathrm F}.\label{eq:fw-convex}
\end{align}
Combining \Cref{eq:fw-comparater} with \Cref{eq:fw-convex} yields
\begin{align*}
 \left\langle
  \nabla L(\Wb_s),\Sb_s-\Wb_s
 \right\rangle_{\mathrm F}
 &\leq
 L(\widehat\Wb_{K_s})-L(\Wb_s).
\end{align*}

Since
$\Wb_{s+1}-\Wb_s=\gamma_s(\Sb_s-\Wb_s)$, the operator-norm smoothness
inequality gives
\begin{align*}
 L(\Wb_{s+1})
 &\leq
 L(\Wb_s)
 +\gamma_s
 \left\langle
  \nabla L(\Wb_s),\Sb_s-\Wb_s
 \right\rangle_{\mathrm F}
 +\frac{\Hop\gamma_s^2}{2}
 \normop{\Sb_s-\Wb_s}^2\\
 &\leq
 L(\Wb_s)
 +\gamma_s\bigl(L(\widehat\Wb_{K_s})-L(\Wb_s)\bigr)
 +2\Hop\gamma_s^2R_s^2\\
 &=(1-\gamma_s)L(\Wb_s)
 +\gamma_sL(\widehat\Wb_{K_s})
 +2\Hop\gamma_s^2R_s^2.
\end{align*}
Substituting $\gamma_s=2/(s+2)$ gives
\begin{align*}
 L(\Wb_{s+1})
 &\leq
 \frac{s}{s+2}L(\Wb_s)
 +\frac{2}{s+2}L(\widehat\Wb_{K_s})
 +\frac{8\Hop}{(s+2)^2}R_s^2.
\end{align*}
Multiplying by $(s+1)(s+2)$ yields
\begin{align*}
 (s+1)(s+2)L(\Wb_{s+1})\leq
 s(s+1)L(\Wb_s)
 +2(s+1)L(\widehat\Wb_{K_s})
 +8\Hop\frac{s+1}{s+2}R_s^2.
\end{align*}
Summing from $s=0$ to $s=t-1$ gives
\begin{align*}
 &\sum_{s=0}^{t-1}
 \Bigl[(s+1)(s+2)L(\Wb_{s+1})-s(s+1)L(\Wb_s)
 \Bigr]\leq
 2\sum_{s=0}^{t-1}(s+1)L(\widehat\Wb_{K_s})
 +8\Hop \sum_{s=0}^{t-1}\frac{s+1}{s+2}R_s^2.
\end{align*}
The sum on the left telescopes, and its initial term is
$0\cdot1\cdot L(\Wb_0)=0$.  Hence
\begin{align*}
 t(t+1)L(\Wb_t)
 &\leq
 2\sum_{s=0}^{t-1}(s+1)L(\widehat\Wb_{K_s})
 +8\Hop
 \sum_{s=0}^{t-1}\frac{s+1}{s+2}R_s^2,
\end{align*}
which is the claimed inequality.
\end{proof}

\begin{lemma}[Two scalar sums]
\label{lem:Muon-scalar-sums}
Let $M\geq1$ and let $t\geq M$ be an integer.  Then
\begin{align*}
 \sum_{s=0}^{t-1}
 (s+1)
 \max\left\{1,\frac{s+1}{M}\right\}^{-(\alpha-1)/\alpha}
 &\leq
 M^{(\alpha-1)/\alpha}t^{1+1/\alpha},\\
 \sum_{s=0}^{t-1}
 \max\left\{1,\frac{s+1}{M}\right\}^{1/\alpha}
 &\leq
 M^{-1/\alpha}t^{1+1/\alpha}.
\end{align*}
\end{lemma}

\begin{proof}
Set $m=s+1$, so that $m=1,\ldots,t$. For the first sum, we claim that for every $1\le m\le t$,
\begin{equation}
m\max\left\{1,\frac{m}{M}\right\}^{-(\alpha-1)/\alpha}
\le
M^{(\alpha-1)/\alpha}m^{1/\alpha}.
\label{eq:first-sum-pointwise}
\end{equation}
Indeed, if $m\le M$, then
\[
m\max\left\{1,\frac{m}{M}\right\}^{-(\alpha-1)/\alpha}
=
m.
\]
Since $\alpha>1$, we have $(\alpha-1)/\alpha>0$, and hence
\[
M^{(\alpha-1)/\alpha}m^{1/\alpha}
=
m\left(\frac{M}{m}\right)^{(\alpha-1)/\alpha}
\ge m.
\]
Thus \eqref{eq:first-sum-pointwise} holds when $m\le M$.

If $m>M$, then
\begin{align*}
m\max\left\{1,\frac{m}{M}\right\}^{-(\alpha-1)/\alpha}
&=
m\left(\frac{m}{M}\right)^{-(\alpha-1)/\alpha}\\
&=
M^{(\alpha-1)/\alpha}
m^{1-(\alpha-1)/\alpha}\\
&=
M^{(\alpha-1)/\alpha}m^{1/\alpha},
\end{align*}
so in this case equality holds in \eqref{eq:first-sum-pointwise}.
Therefore, using $m\le t$,
\begin{align*}
\sum_{s=0}^{t-1}
(s+1)
\max\left\{1,\frac{s+1}{M}\right\}^{-(\alpha-1)/\alpha}
&=
\sum_{m=1}^{t}
m\max\left\{1,\frac{m}{M}\right\}^{-(\alpha-1)/\alpha}\\
&\le
M^{(\alpha-1)/\alpha}
\sum_{m=1}^{t}m^{1/\alpha}\\
&\le
M^{(\alpha-1)/\alpha}
t \cdot t^{1/\alpha}=
M^{(\alpha-1)/\alpha}t^{1+1/\alpha}.
\end{align*}

For the second sum, since $t\ge M$, we have $\frac{t}{M}\ge1.$ We again verify the required bound pointwise.  If $m\le M$, then
\[
\max\left\{1,\frac{m}{M}\right\}^{1/\alpha}
=
1
\le
\left(\frac{t}{M}\right)^{1/\alpha}.
\]
If $M<m\le t$, then
\[
\max\left\{1,\frac{m}{M}\right\}^{1/\alpha}
=
\left(\frac{m}{M}\right)^{1/\alpha}
\le
\left(\frac{t}{M}\right)^{1/\alpha}.
\]
Hence, for every $1\le m\le t$,
\[
\max\left\{1,\frac{m}{M}\right\}^{1/\alpha}
\le
\left(\frac{t}{M}\right)^{1/\alpha}.
\]
Summing over $m=1,\ldots,t$ gives
\begin{align*}
\sum_{s=0}^{t-1}
\max\left\{1,\frac{s+1}{M}\right\}^{1/\alpha}
&=
\sum_{m=1}^{t}
\max\left\{1,\frac{m}{M}\right\}^{1/\alpha}\\
&\le
t\left(\frac{t}{M}\right)^{1/\alpha}= M^{-1/\alpha}t^{1+1/\alpha}.
\end{align*}
This proves both bounds.
\end{proof}

\begin{proof}[Proof of \Cref{thm:Muon-upper bound}]
Work on the intersection of the event in \Cref{prop:prefix-v2} and the
regularity event in \Cref{lem:regularity-v2}.  Its probability is at least
$1-O(d^{-q})$, and \Cref{lem:smoothness-v2} gives $\Hop\leq\frac{1221}{2000}.$
Let $A_{\alpha,q}$ and $c_*$ be the constants in \Cref{lem:Muon-time-comparators}.  Choose
\begin{align*}
 C_{\mathrm{low}}
 &\geq1,
 \qquad
 c_{\mathrm{up}}\leq c_*.
\end{align*}
Fix an integer $t$ in the range of the theorem.
The lower time bound and $C_{\mathrm{low}}\geq1$ give $t\geq M$.

The radius sequence is nondecreasing.  For every $0\leq s<t$, write the
oracle point in the theorem as
\begin{align*}
 \Sb_s
 &:=-R_s\polar\!\left(\nabla L(\Wb_s^{\mathrm{Muon}})\right).
\end{align*}
The defining properties of the polar factor and operator--nuclear norm
duality give
\begin{align*}
 \normop{\Sb_s}
 &\leq R_s,\\
 \left\langle
  \nabla L(\Wb_s^{\mathrm{Muon}}),\Sb_s
 \right\rangle_{\mathrm F}
 &=-R_s\left\|\nabla L(\Wb_s^{\mathrm{Muon}})\right\|_*\\
 &=
 \min_{\normop{\Sb}\leq R_s}
 \left\langle
  \nabla L(\Wb_s^{\mathrm{Muon}}),\Sb
 \right\rangle_{\mathrm F}.
\end{align*}
Moreover, the recursion in the theorem is exactly
\begin{align*}
 \Wb_{s+1}^{\mathrm{Muon}}
 &=(1-\gamma_s)\Wb_s^{\mathrm{Muon}}+\gamma_s\Sb_s.
\end{align*}

For every $0\leq s<t$, use the prefix comparator
$\widehat\Wb_{K_s}$ from \Cref{lem:Muon-time-comparators}.  The first
conclusion of that lemma gives
\begin{align*}
 \normop{\widehat\Wb_{K_s}}\leq
 C_\alpha(\log N)
 \max\left\{1,\left(\frac{s+1}{M}\right)^{1/(2\alpha)}\right\}=R_s.
\end{align*}
Therefore \Cref{lem:Muon-changing-radius} yields
\begin{align*}
 t(t+1)L(\Wb_t^{\mathrm{Muon}})
 &\leq
 2\sum_{s=0}^{t-1}(s+1)L(\widehat\Wb_{K_s})
 +8\Hop
 \sum_{s=0}^{t-1}\frac{s+1}{s+2}R_s^2.
\end{align*}

We bound the comparator contribution first.  The second conclusion of
\Cref{lem:Muon-time-comparators} and the first estimate in
\Cref{lem:Muon-scalar-sums} give
\begin{align*}
 \sum_{s=0}^{t-1}(s+1)L(\widehat\Wb_{K_s})
 &\leq
 A_{\alpha,q}(\log N)d^{1-\alpha}
 \sum_{s=0}^{t-1}
 (s+1)
 \max\left\{1,\frac{s+1}{M}\right\}^{-(\alpha-1)/\alpha}\\
 &\leq
 A_{\alpha,q}(\log N)d^{1-\alpha}
 M^{(\alpha-1)/\alpha}t^{1+1/\alpha}.
\end{align*}

For the curvature contribution, the definition of $R_s$ gives
\begin{align*}
 R_s^2
 &=C_\alpha^2(\log N)^2
 \max\left\{1,\frac{s+1}{M}\right\}^{1/\alpha}.
\end{align*}
Using $(s+1)/(s+2)\leq1$ and the second estimate in
\Cref{lem:Muon-scalar-sums}, we obtain
\begin{align*}
 \sum_{s=0}^{t-1}\frac{s+1}{s+2}R_s^2
 &\leq
 C_\alpha^2(\log N)^2
 \sum_{s=0}^{t-1}
 \max\left\{1,\frac{s+1}{M}\right\}^{1/\alpha}\\
 &\leq
 C_\alpha^2(\log N)^2
 M^{-1/\alpha}t^{1+1/\alpha}.
\end{align*}

Substituting the last two estimates into the changing-radius inequality and
using $t(t+1)\geq t^2$ gives
\begin{align*}
 L(\Wb_t^{\mathrm{Muon}})
 &\leq
 2A_{\alpha,q}(\log N)d^{1-\alpha}
 M^{(\alpha-1)/\alpha}t^{-(\alpha-1)/\alpha}\\
 &\qquad+
 8\Hop C_\alpha^2(\log N)^2
 M^{-1/\alpha}t^{-(\alpha-1)/\alpha}.
\end{align*}
The two coefficients have exactly the same scale.  Indeed, since
$M=(\log N)d^{\alpha-1}$,
\begin{align*}
 \frac{(\log N)^2M^{-1/\alpha}}
 {(\log N)d^{1-\alpha}M^{(\alpha-1)/\alpha}}
 &=(\log N)d^{\alpha-1}
 M^{-1/\alpha-(\alpha-1)/\alpha}\\
 &=M\cdot M^{-1}=1.
\end{align*}
Consequently, the bound $\Hop\leq1221/2000$ gives
\begin{align*}
 L(\Wb_t^{\mathrm{Muon}})
 &\leq
 \left(
  2A_{\alpha,q}
  +\frac{1221}{250}C_\alpha^2
 \right)
 (\log N)d^{1-\alpha}
 M^{(\alpha-1)/\alpha}t^{-(\alpha-1)/\alpha}.
\end{align*}
Finally, substituting $M=(\log N)d^{\alpha-1}$ and simplifying each power
gives
\begin{align*}
 (\log N)d^{1-\alpha}
 M^{(\alpha-1)/\alpha}t^{-(\alpha-1)/\alpha}=
 (\log N)\left(\frac{dt}{\log N}\right)^{-(\alpha-1)/\alpha}.
\end{align*}
This proves the asserted loss bound for the fixed $t$.  Since \Cref{prop:prefix-v2} holds on a single event simultaneously for all
admissible $K$, all comparator bounds used above hold simultaneously for
every $K_s$. The remaining argument is deterministic, so the resulting
loss bound holds on the same event for every integer $t$ in the stated range.
\end{proof}

% ============================================================



% ============================================================
\section{Proof of \Cref{thm:gd}}
% ============================================================

\subsection{Block gradients and nonexpansiveness}

For $B\subset[N]$, write
$L_B(\Wb):=\sum_{i\in B}p_i\ell_i(\Wb)$ only in this section.

\begin{lemma}[Uniform block-gradient bound]
\label{lem:block-v2}
On the event in \Cref{lem:regularity-v2}, every consecutive block with
$d\leq|B|\leq2d$ and $\min B\geq d/3$ satisfies
\[
 \sup_{\Wb}\normF{\nabla L_B(\Wb)}^2
 \lesssim\sum_{i\in B}p_i^2.
\]
\end{lemma}

\begin{proof}
Differentiating the loss once gives
\[
 \nabla\ell_i(\Wb)
 =\left(\sum_{j=1}^N\pi_{ji}(\Wb)\ub_j-\ub_i\right)\vb_i^\top.
\]
Fix $\Ab$ with $\normF{\Ab}=1$.  Expanding the Frobenius inner product,
then applying the triangle inequality and Euclidean Cauchy--Schwarz, gives
\begin{align*}
 |\langle\Ab,\nabla L_B(\Wb)\rangle_{\mathrm F}|&=
 \left|
  \sum_{i\in B}p_i
  \left\langle
   \Ab,
   \left(\sum_j\pi_{ji}(\Wb)\ub_j-\ub_i\right)\vb_i^\top
  \right\rangle_{\mathrm F}
 \right|\\
 &\quad=
 \left|
  \sum_{i\in B}p_i
  \left(\sum_j\pi_{ji}(\Wb)\ub_j-\ub_i\right)^\top
  \Ab\vb_i
 \right|\\
 &\quad\leq
 \sum_{i\in B}p_i
 \left\|\sum_j\pi_{ji}(\Wb)\ub_j-\ub_i\right\|
 \norm{\Ab\vb_i}\\
 &\quad\leq3\sum_{i\in B}p_i\norm{\Ab\vb_i}.
\end{align*}
The last line uses
\begin{align*}
 \left\|\sum_j\pi_{ji}(\Wb)\ub_j-\ub_i\right\|
 &\leq\sum_j\pi_{ji}(\Wb)\norm{\ub_j}+\norm{\ub_i}\\
 &\leq2\sqrt{11/10}<3.
\end{align*}
Weighted scalar Cauchy--Schwarz now gives
\begin{align*}
 \sum_{i\in B}p_i\norm{\Ab\vb_i}
 &=\sum_{i\in B}\sqrt{p_i}
 \bigl(\sqrt{p_i}\norm{\Ab\vb_i}\bigr)\\
 &\leq
 \left(\sum_{i\in B}p_i\right)^{1/2}
 \left(\sum_{i\in B}p_i\norm{\Ab\vb_i}^2\right)^{1/2}.
\end{align*}

We next expand the matrix term that appears here.  Since
$\normF{\Ab}=1$,
\begin{align*}
 \sum_{i\in B}p_i\norm{\Ab\vb_i}^2
 &=\sum_{i\in B}p_i\vb_i^\top\Ab^\top\Ab\vb_i\\
 &=\operatorname{tr}\!\left(
  \Ab^\top\Ab\sum_{i\in B}p_i\vb_i\vb_i^\top
 \right)\\
 &\leq
 \lambda_{\max}\!\left(\sum_{i\in B}p_i\vb_i\vb_i^\top\right)
 \operatorname{tr}(\Ab^\top\Ab)\\
 &=\lambda_{\max}\!\left(\sum_{i\in B}p_i\vb_i\vb_i^\top\right).
\end{align*}
The inequality follows because for a positive-semidefinite matrix $M$,
$M\preceq\lambda_{\max}(M)\Ib_d$, and multiplication by
$\Ab^\top\Ab\succeq0$ followed by the trace preserves this order.

For every unit vector $\zb$,
\begin{align*}
 \zb^\top\left(\sum_{i\in B}p_i\vb_i\vb_i^\top\right)\zb
 &=\sum_{i\in B}p_i(\vb_i^\top\zb)^2\\
 &\leq\left(\max_{i\in B}p_i\right)
 \sum_{i\in B}(\vb_i^\top\zb)^2\\
 &=\left(\max_{i\in B}p_i\right)
 \left\|[\vb_i]_{i\in B}^\top\zb\right\|^2\\
 &\leq\left(\max_{i\in B}p_i\right)
 \normop{[\vb_i]_{i\in B}}^2\\
 &\leq16\max_{i\in B}p_i.
\end{align*}
Taking the supremum over unit $\zb$ proves the corresponding eigenvalue
bound.

Finally, since $B$ is consecutive, $|B|\le 2d$, and $\min B\ge d/3$,
\[
\frac{\max B}{\min B}
\le
1+\frac{2d}{\min B}
\le 7.
\]
The largest index in $B$ is at most seven times the smallest, so
$\max_{i\in B}p_i\leq7^\alpha\min_{i\in B}p_i$.  Hence
\begin{align*}
 \left(\sum_{i\in B}p_i\right)\max_{i\in B}p_i
 &\leq7^\alpha
 \left(\sum_{i\in B}p_i\right)\min_{i\in B}p_i\\
 &\leq7^\alpha\sum_{i\in B}p_i^2,
\end{align*}
where the last inequality uses
$p_i^2\geq(\min_{k\in B}p_k)p_i$ term by term.  Substitution into all
preceding displays gives
\[
 |\langle\Ab,\nabla L_B(\Wb)\rangle_{\mathrm F}|
 \lesssim\left(\sum_{i\in B}p_i^2\right)^{1/2}.
\]
Taking the supremum over all Frobenius-unit $\Ab$ and then squaring proves
the lemma.
\end{proof}

\begin{lemma}[Cocoercivity and nonexpansiveness]
\label{lem:cocoercive-v2}
If $f$ is convex and $H$-smooth in Frobenius norm, then
\[
 \left\langle
  \nabla f(\Xb)-\nabla f(\Yb),\Xb-\Yb
 \right\rangle_{\mathrm F}
 \geq\frac1H\normF{\nabla f(\Xb)-\nabla f(\Yb)}^2.
\]
Consequently, $\Wb\mapsto\Wb-\eta\nabla f(\Wb)$ is nonexpansive whenever
$0\leq\eta\leq2/H$.
\end{lemma}

\begin{proof}
Fix $\Xb,\Yb$ and define
$\phi(\Zb):=f(\Zb)-\langle\nabla f(\Yb),\Zb\rangle_{\mathrm F}$.
Then $\nabla\phi(\Yb)=\mathbf0$, so convexity makes $\Yb$ a minimizer of
$\phi$.  Smoothness at $\Xb$ in direction
$-H^{-1}\nabla\phi(\Xb)$ gives
\begin{align*}
 \phi\!\left(\Xb-\frac1H\nabla\phi(\Xb)\right)
 &\leq\phi(\Xb)
 -\frac1H\normF{\nabla\phi(\Xb)}^2
 +\frac{H}{2}\frac1{H^2}\normF{\nabla\phi(\Xb)}^2\\
 &=\phi(\Xb)-\frac1{2H}\normF{\nabla\phi(\Xb)}^2.
\end{align*}
Since $\phi(\Yb)$ is no larger than the left-hand side, substitution of the
definition of $\phi$ gives
\[
 f(\Xb)-f(\Yb)
 -\langle\nabla f(\Yb),\Xb-\Yb\rangle_{\mathrm F}
 \geq\frac1{2H}\normF{\nabla f(\Xb)-\nabla f(\Yb)}^2.
\]
Interchanging $\Xb$ and $\Yb$ gives the second inequality
\[
 f(\Yb)-f(\Xb)
 -\langle\nabla f(\Xb),\Yb-\Xb\rangle_{\mathrm F}
 \geq\frac1{2H}\normF{\nabla f(\Xb)-\nabla f(\Yb)}^2.
\]
Adding their left-hand sides cancels the function values and yields
\begin{align*}
 &-\langle\nabla f(\Yb),\Xb-\Yb\rangle_{\mathrm F}
 -\langle\nabla f(\Xb),\Yb-\Xb\rangle_{\mathrm F}\\
 &\qquad=
 \langle\nabla f(\Xb)-\nabla f(\Yb),\Xb-\Yb\rangle_{\mathrm F}.
\end{align*}
Adding the right-hand sides proves cocoercivity.

Put $\Gb:=\nabla f(\Xb)-\nabla f(\Yb)$.  Direct expansion gives
\begin{align*}
 &\left\|
  (\Xb-\eta\nabla f(\Xb))-(\Yb-\eta\nabla f(\Yb))
 \right\|_{\mathrm F}^2\\
 &\qquad=
 \normF{\Xb-\Yb}^2
 -2\eta\langle\Gb,\Xb-\Yb\rangle_{\mathrm F}
 +\eta^2\normF{\Gb}^2\\
 &\qquad\leq
 \normF{\Xb-\Yb}^2
 -\left(\frac{2\eta}{H}-\eta^2\right)\normF{\Gb}^2\\
 &\qquad\leq\normF{\Xb-\Yb}^2
\end{align*}
when $0\leq\eta\leq2/H$.
\end{proof}

Fix a block $B$ and run the same schedule after deleting that block:
\[
 \widetilde\Wb_0^{(B)}=\mathbf0,
 \qquad
 \widetilde\Wb_{t+1}^{(B)}
 =\widetilde\Wb_t^{(B)}
 -\eta_t\nabla(L-L_B)(\widetilde\Wb_t^{(B)}).
\]

\begin{lemma}[Leave-block stability]
\label{lem:leave-v2}
On the regularity event, every block in \Cref{lem:block-v2} satisfies
\begin{align*}
 \normF{\Wb_t^{\mathrm{GD}}-\widetilde\Wb_t^{(B)}}
 &\lesssim
 \left(\sum_{s<t}\eta_s\right)
 \left(\sum_{i\in B}p_i^2\right)^{1/2},\\
 L_B(\Wb_t^{\mathrm{GD}})
 &\geq L_B(\widetilde\Wb_t^{(B)})
 -O\!\left(
  \left(\sum_{s<t}\eta_s\right)
  \sum_{i\in B}p_i^2
 \right).
\end{align*}
\end{lemma}

\begin{proof}
Let
$T_t(\Wb):=\Wb-\eta_t\nabla L(\Wb)$.  By
\Cref{lem:smoothness-v2,lem:cocoercive-v2}, $T_t$ is nonexpansive.  The
reduced update is exactly
\[
 \widetilde\Wb_{t+1}^{(B)}
 =T_t(\widetilde\Wb_t^{(B)})
 +\eta_t\nabla L_B(\widetilde\Wb_t^{(B)}).
\]
Subtracting it from the full update gives
\begin{align*}
 \Wb_{t+1}^{\mathrm{GD}}-\widetilde\Wb_{t+1}^{(B)}
 &=T_t(\Wb_t^{\mathrm{GD}})-T_t(\widetilde\Wb_t^{(B)})-\eta_t\nabla L_B(\widetilde\Wb_t^{(B)}).
\end{align*}
Taking Frobenius norms, using nonexpansiveness, and then applying
\Cref{lem:block-v2} gives
\begin{align*}
 \normF{\Wb_{t+1}^{\mathrm{GD}}-\widetilde\Wb_{t+1}^{(B)}}
 &\leq\normF{\Wb_t^{\mathrm{GD}}-\widetilde\Wb_t^{(B)}}+\eta_t
 \sup_{\Zb}\normF{\nabla L_B(\Zb)}\\
 &\leq\normF{\Wb_t^{\mathrm{GD}}-\widetilde\Wb_t^{(B)}}
 +O\!\left(
  \eta_t\left(\sum_{i\in B}p_i^2\right)^{1/2}
 \right).
\end{align*}
The initial distance is zero.  Summing this scalar recurrence from
$s=0$ to $s=t-1$ proves the first claim.

For arbitrary $\Wb,\widetilde\Wb$, the fundamental theorem of calculus
along the line segment gives
\begin{align*}
 L_B(\Wb)-L_B(\widetilde\Wb)
 &=\int_0^1
 \left\langle
  \nabla L_B(\widetilde\Wb+s(\Wb-\widetilde\Wb)),
  \Wb-\widetilde\Wb
 \right\rangle_{\mathrm F}\,ds\\
 &\geq-\int_0^1
 \normF{\nabla L_B(\widetilde\Wb+s(\Wb-\widetilde\Wb))}
 \normF{\Wb-\widetilde\Wb}\,ds\\
 &\geq-\sup_{\Zb}\normF{\nabla L_B(\Zb)}
 \normF{\Wb-\widetilde\Wb}.
\end{align*}
Apply this inequality with the two trajectories, then substitute the first
claim and \Cref{lem:block-v2}.  This proves the second claim.
\end{proof}

\subsection{Sign symmetry and preservation of the tail}

\begin{lemma}[Sign-symmetry bound]
\label{lem:sign-v2}
Fix the keys, vectors $(\zb_i)_{i\in B}$, and a matrix $\Wb$.  Put
$\vb_i=\sigma_i\zb_i$ for independent Rademacher signs.  For every
$0<\rho<1$,
\[
 \PP_\sigma\!\left(
  L_B(\Wb)<(1-\rho)(\log N)\sum_{i\in B}p_i
 \right)
 \leq
 \exp\!\left(
  -\frac{\rho^2(\sum_{i\in B}p_i)^2}
  {2\sum_{i\in B}p_i^2}
 \right).
\]
\end{lemma}

\begin{proof}
Display the query as $\ell_i(\Wb;\vb)$.  It is convex in $\vb$ and equals
$\log N$ at zero, so midpoint convexity gives
\[
 \ell_i(\Wb;\zb_i)+\ell_i(\Wb;-\zb_i)\geq2\log N.
\]
Let
$Y_i(\sigma_i):=\min\{\ell_i(\Wb;\sigma_i\zb_i),2\log N\}$.  Since the
losses are nonnegative, either truncation changes neither signed loss, or
one truncated loss alone equals $2\log N$.  Therefore
\[
 \frac{Y_i(+1)+Y_i(-1)}2\geq\log N.
\]
The independent variables $p_iY_i(\sigma_i)$ lie in
$[0,2p_i\log N]$.  If
$Z:=\sum_{i\in B}p_iY_i(\sigma_i)$ and
$b:=(\log N)\sum_{i\in B}p_i$, then $\EE_\sigma Z\geq b$.  Hence
\begin{align*}
 \{Z<(1-\rho)b\}
 &\subseteq
 \{Z-\EE_\sigma Z\leq-\rho b\}.
\end{align*}
Hoeffding's inequality, with interval lengths $2p_i\log N$, gives
\begin{align*}
 \PP_\sigma(Z<(1-\rho)b)
 &\leq
 \exp\!\left(
  -\frac{2\rho^2b^2}{\sum_{i\in B}(2p_i\log N)^2}
 \right)\\
 &=\exp\!\left(
  -\frac{\rho^2(\sum_{i\in B}p_i)^2}
  {2\sum_{i\in B}p_i^2}
 \right).
\end{align*}
Finally,
$L_B(\Wb)\geq\sum_{i\in B}p_iY_i(\sigma_i)=Z$, so the same bound applies
to the event in the lemma.
\end{proof}

\begin{proposition}[Tail preservation]
\label{prop:tail-v2}
Fix the target, schedule, a constant $0<\rho<1$, an integer
$d/3\leq r\leq d^2/\log N$, and
$0\leq T\leq\lceil d^{2\alpha+2}\rceil$.  With probability at least
$1-e^{-\Omega(d)}$, simultaneously for $0\leq t\leq T$,
\begin{align*}
 \sum_{i=r}^Np_i\ell_i(\Wb_t^{\mathrm{GD}})
 &\geq(1-\rho)(\log N)\sum_{i=r}^Np_i-O\!\left(
  \left(\sum_{s<t}\eta_s\right)
  \sum_{i=r}^Np_i^2
 \right).
\end{align*}
\end{proposition}

\begin{proof}
Let $m=N-r+1$ and $J=\lfloor m/d\rfloor$.  Since
$r\leq d^2/\log N$ and $N\gtrsim d^2$, one has $m\geq2d$ for large $d$.  Divide
$\{r,\ldots,N\}$ into $J$ consecutive blocks whose sizes differ by at most
one.  The inequalities
\[
 J\leq\frac md,
 \qquad
 J\geq\frac md-1\geq\frac m{2d}
\]
show that every block has size between $d$ and $2d$.  Its first index is at
least $r\geq d/3$.

Represent each query vector as
$\vb_i=\sigma_i\zb_i,$
where $\zb_i\sim\mathcal N(0,\Ib_d/d)$ are independent and
$\sigma_i\in\{-1,+1\}$ are independent Rademacher variables, independent
also of the $\zb_i$.  Since the Gaussian distribution is symmetric,
$\sigma_i\zb_i\stackrel{d}{=}\zb_i$ for every $i$, and independence is
preserved across indices.  Hence the family
$(\vb_i)_{i=1}^N=(\sigma_i\zb_i)_{i=1}^N$ has exactly the same joint law
as the original Gaussian queries.

Fix a block $B$ and a time $t$.  Let $\mathcal C_B$ be the
$\sigma$-algebra generated by all key vectors $(\ub_i)_{i=1}^N$, all
Gaussian magnitudes/directions $(\zb_i)_{i=1}^N$, and all signs
$(\sigma_i)_{i\notin B}$.  Conditional on $\mathcal C_B$, the only
remaining randomness is therefore the collection
$(\sigma_i)_{i\in B}$, which remains a family of independent Rademacher
variables.

Recall that the leave-block trajectory is defined by
\[
 \widetilde\Wb^{(B)}_{s+1}
 =
 \widetilde\Wb^{(B)}_s
 -
 \eta_s\,
 \nabla(L-L_B)(\widetilde\Wb^{(B)}_s),
 \qquad
 \widetilde\Wb^{(B)}_0=\mathbf 0.
\]
The objective $L-L_B$ contains no loss term indexed by $i\in B$ and hence
does not involve the corresponding query vectors
$\vb_i=\sigma_i\zb_i$, $i\in B$.  Since the schedule
$(\eta_s)_{s\ge0}$ is prescribed before the sample is drawn, it is
also independent of these signs.  Consequently,
$\widetilde\Wb^{(B)}_s$ is $\mathcal C_B$-measurable for every $s$,
and in particular $\widetilde\Wb^{(B)}_t$ is fixed when we condition on
$\mathcal C_B$.  We may therefore apply the sign-symmetry bound to
$\widetilde\Wb^{(B)}_t$ using the still-random signs
$(\sigma_i)_{i\in B}$.

The regularity event $\mathcal E_{\mathrm{reg}}$ from \Cref{lem:regularity-v2} is also
$\mathcal C_B$-measurable: query norms and absolute inner products are
sign-invariant, and multiplying columns by signs preserves singular values.
Weights within $B$ differ by at most $7^\alpha$, so
\begin{align*}
 \frac{(\sum_{i\in B}p_i)^2}{\sum_{i\in B}p_i^2}\geq
 \frac{|B|^2(\min_{i\in B}p_i)^2}
 {|B|(\max_{i\in B}p_i)^2}\geq\frac{|B|}{7^{2\alpha}}
 \geq\frac d{7^{2\alpha}}.
\end{align*}
Thus \Cref{lem:leave-v2,lem:sign-v2} imply, except with conditional
probability $e^{-c_{\alpha,\rho}d}$,
\begin{align*}
 L_B(\Wb_t^{\mathrm{GD}})
 &\geq L_B(\widetilde\Wb_t^{(B)})
 -O\!\left(
  \left(\sum_{s<t}\eta_s\right)
  \sum_{i\in B}p_i^2
 \right)\\
 &\geq(1-\rho)(\log N)\sum_{i\in B}p_i
 -O\!\left(
  \left(\sum_{s<t}\eta_s\right)
  \sum_{i\in B}p_i^2
 \right).
\end{align*}

Let $\mathcal F_{B,t}$ be failure of this block inequality.  The tower
property and $\mathcal C_B$-measurability of the regularity event give
\begin{align*}
 \PP(\mathcal E_{\mathrm{reg}}\cap\mathcal F_{B,t})
 &=\EE\!\left[
  \EE[\mathbf{1}_{\mathcal E_{\mathrm{reg}}}
      \mathbf{1}_{\mathcal F_{B,t}}\mid\mathcal C_B]
 \right]\\
 &=\EE\!\left[
  \mathbf{1}_{\mathcal E_{\mathrm{reg}}}
  \PP(\mathcal F_{B,t}\mid\mathcal C_B)
 \right]\\
 &\leq e^{-c_{\alpha,\rho}d}.
\end{align*}
There are at most $N/d$ blocks and polynomially many times.  A union bound,
followed by adding
$\PP(\mathcal E_{\mathrm{reg}}^c)\leq e^{-\Omega(d)}$, keeps total failure
at $e^{-\Omega(d)}$.

On the complementary event, summing over the exact block partition gives
\begin{align*}
 \sum_B L_B(\Wb_t^{\mathrm{GD}})
 &=\sum_{i=r}^Np_i\ell_i(\Wb_t^{\mathrm{GD}}),\\
 \sum_B\sum_{i\in B}p_i
 &=\sum_{i=r}^Np_i,\\
 \sum_B\sum_{i\in B}p_i^2
 &=\sum_{i=r}^Np_i^2.
\end{align*}
Substitution proves the proposition.
\end{proof}

% \subsection{Completion of the GD lower bound}

% \begin{proof}[Proof of \Cref{thm:gd}]
% Choose an integer
% $ K\eqsim\left(\frac{\log N}{\eps}\right)^{1/(\alpha-1)}
% $
% with a sufficiently small fixed implicit constant.  The target range makes
% $d\leq K\leq d^2/\log N$, and the implicit constant can be chosen so that
% $Kp_K\log N$ is a sufficiently large fixed multiple of $\eps$.  Put
% $r=\lfloor K/2\rfloor+1$.

% We record the three weight estimates with their derivations.  There are at
% least $K/3$ indices in $\{r,\ldots,K\}$ for large $K$, and each has weight
% at least $p_K$, so
% \[
%  \sum_{i=r}^Kp_i\geq\frac K3p_K\gtrsim Kp_K.
% \]
% Since $r>K/2$, integral comparison gives
% \begin{align*}
%  \sum_{i=r}^Np_i
%  &\leq
%  \frac{r^{-\alpha}+\int_r^\infty y^{-\alpha}\,dy}
%  {\sum_{j=1}^Nj^{-\alpha}}\lesssim\frac{K^{1-\alpha}}{\sum_{j=1}^Nj^{-\alpha}}
%  \eqsim Kp_K,\\
%  \sum_{i=r}^Np_i^2
%  &\leq
%  \frac{r^{-2\alpha}+\int_r^\infty y^{-2\alpha}\,dy}
%  {(\sum_{j=1}^Nj^{-\alpha})^2}\lesssim\frac{K^{1-2\alpha}}
%  {(\sum_{j=1}^Nj^{-\alpha})^2}
%  \eqsim Kp_K^2.
% \end{align*}

% Apply \Cref{prop:tail-v2} with a fixed sufficiently small $\rho>0$.
% The weight bounds above imply, after choosing this fixed $\rho$ sufficiently
% small, that
% \begin{align*}
%  &(\log N)\sum_{i=r}^Kp_i
%  -\rho(\log N)\sum_{i=r}^Np_i \gtrsim Kp_K\log N.
% \end{align*}
% Using the exact decomposition
% $\sum_{i=r}^Np_i=\sum_{i>K}p_i+\sum_{i=r}^Kp_i$, we obtain
% \begin{align*}
%  L(\Wb_t^{\mathrm{GD}})
%  &\geq\sum_{i=r}^Np_i\ell_i(\Wb_t^{\mathrm{GD}})\\
%  &\geq(1-\rho)(\log N)\sum_{i=r}^Np_i
%  -O\!\left(
%   \left(\sum_{s<t}\eta_s\right)
%   \sum_{i=r}^Np_i^2
%  \right)\\
%  &= (\log N)\sum_{i>K}p_i
%  +(\log N)\sum_{i=r}^Kp_i
%  -\rho(\log N)\sum_{i=r}^Np_i-O\!\left(
%   \left(\sum_{s<t}\eta_s\right)
%   \sum_{i=r}^Np_i^2
%  \right)\\
%  &\geq(\log N)\sum_{i>K}p_i
%  +\left[
%   (\log N)\sum_{i=r}^Kp_i
%   -\rho(\log N)\sum_{i=r}^Np_i
%  \right]-
%  O\!\left(
%   \left(\sum_{s<t}\eta_s\right)Kp_K^2
%  \right).
% \end{align*}
% If
% \[
%  \sum_{s<t}\eta_s
%  \lesssim\frac{\log N}{p_K},
% \]
% where the implicit constant on the left is chosen sufficiently small, the
% last error term is at most one half of the positive bracket.  Substitution
% into the last display therefore gives
% \begin{align*}
%  L(\Wb_t^{\mathrm{GD}})
%  &\geq(\log N)\sum_{i>K}p_i
%  +\frac12\left[
%   (\log N)\sum_{i=r}^Kp_i
%   -\rho(\log N)\sum_{i=r}^Np_i
%  \right]\\
%  &\gtrsim Kp_K\log N >\eps,
% \end{align*}
% where the strict last inequality is ensured by the fixed implicit constant
% in the choice of $K$.

% Because $\eta_s\leq3$,
% \[
%  \sum_{s<t}\eta_s\leq3t.
% \]
% Consequently, every $t\lesssim(\log N)/p_K$, with a sufficiently small
% implicit constant, lies in the preceding cumulative-step regime and has
% loss larger than the target.  This time is within the horizon in
% \Cref{prop:tail-v2}, because
% \begin{align*}
%  \frac{\log N}{p_K}
%  &\lesssim(\log N)K^\alpha\lesssim(\log N)\left(\frac{d^2}{\log N}\right)^\alpha
%  \leq d^{2\alpha+2}
% \end{align*}
% for large $d$.  Hence
% \begin{align*}
%  \tau_{\mathrm{GD}}(\eps)
%  &\gtrsim\frac{\log N}{p_K}\eqsim(\log N)K^\alpha\eqsim(\log N)
%  \left(\frac{\log N}{\eps}\right)^{\alpha/(\alpha-1)}.
% \end{align*}
% Intersecting this pointwise schedule event with the common event in
% \Cref{thm:fw} and applying the Muon upper bound there proves
% $\tau_{\mathrm{GD}}(\eps)\gtrsim d\tau_{\mathrm{Muon}}(\eps)$.
% The probability statement remains pointwise in the fixed target--schedule
% pair because the leave-block trajectory is independent of the removed signs
% only for a schedule fixed before the sample.
% \end{proof}


\begin{proof}[Proof of \Cref{thm:gd}]
Fix an arbitrary target $\eps$ in the stated range and an arbitrary deterministic pre-sample schedule satisfying $\eta_t<3.$ We first prove a $K$-parameterized lower bound that holds for every integer $d\leq K\leq\frac{d^2}{\log N},$ without yet relating $K$ to $\eps$. For each such $K$,  Put $r=\lfloor K/2\rfloor+1$, which is admissible in \Cref{prop:tail-v2}.

We record the three weight estimates with their derivations.  There are at
least $K/3$ indices in $\{r,\ldots,K\}$ for large $K$, and each has weight
at least $p_K$, so
\[
 \sum_{i=r}^Kp_i\geq\frac K3p_K\gtrsim Kp_K.
\]
Since $r>K/2$, integral comparison gives
\begin{align*}
 \sum_{i=r}^Np_i
 &\leq
 \frac{r^{-\alpha}+\int_r^\infty y^{-\alpha}\,dy}
 {\sum_{j=1}^Nj^{-\alpha}}\lesssim\frac{K^{1-\alpha}}{\sum_{j=1}^Nj^{-\alpha}}
 \eqsim Kp_K,\\
 \sum_{i=r}^Np_i^2
 &\leq
 \frac{r^{-2\alpha}+\int_r^\infty y^{-2\alpha}\,dy}
 {(\sum_{j=1}^Nj^{-\alpha})^2}\lesssim\frac{K^{1-2\alpha}}
 {(\sum_{j=1}^Nj^{-\alpha})^2}
 \eqsim Kp_K^2.
\end{align*}

Apply \Cref{prop:tail-v2} with a fixed sufficiently small $\rho>0$.
The weight bounds above imply, after choosing this fixed $\rho$ sufficiently
small, that
\begin{align*}
 &(\log N)\sum_{i=r}^Kp_i
 -\rho(\log N)\sum_{i=r}^Np_i \gtrsim Kp_K\log N.
\end{align*}

Using the exact decomposition
$\sum_{i=r}^Np_i=\sum_{i>K}p_i+\sum_{i=r}^Kp_i$, we obtain
\begin{align*}
 L(\Wb_t^{\mathrm{GD}})
 &\geq\sum_{i=r}^Np_i\ell_i(\Wb_t^{\mathrm{GD}})\\
 &\geq(1-\rho)(\log N)\sum_{i=r}^Np_i
 -O\!\left(
  \left(\sum_{s<t}\eta_s\right)
  \sum_{i=r}^Np_i^2
 \right)\\
 &= (\log N)\sum_{i>K}p_i
 +(\log N)\sum_{i=r}^Kp_i
 -\rho(\log N)\sum_{i=r}^Np_i-O\!\left(
  \left(\sum_{s<t}\eta_s\right)
  \sum_{i=r}^Np_i^2
 \right)\\
 &\geq(\log N)\sum_{i>K}p_i
 +\left[
  (\log N)\sum_{i=r}^Kp_i
  -\rho(\log N)\sum_{i=r}^Np_i
 \right]-
 O\!\left(
  \left(\sum_{s<t}\eta_s\right)Kp_K^2
 \right).
\end{align*}

Consequently, if
$ \sum_{s<t}\eta_s
 \lesssim\frac{\log N}{p_K},
$
where the implicit constant is chosen sufficiently small, the
last error term is at most one half of the positive bracket.  Substitution
into the last display therefore gives
\begin{align*}
 L(\Wb_t^{\mathrm{GD}})
 &\geq(\log N)\sum_{i>K}p_i
 +\frac12\left[
  (\log N)\sum_{i=r}^Kp_i
  -\rho(\log N)\sum_{i=r}^Np_i
 \right]\\
 &\gtrsim Kp_K\log N,
\end{align*}

In the application of \Cref{prop:tail-v2}, take $T=\lceil d^{2\alpha+2}\rceil$.  Since there are at most $d^2$ admissible
integer values of $K$, a union bound shows that, with probability at least $1-e^{-\Omega(d)},$ the preceding estimate holds simultaneously for every admissible $K$ and
every $0\leq t\leq T$.

Since $0\leq\eta_s\leq3$, whenever $t$ is at most a sufficiently
small fixed multiple of $(\log N)/p_K$, we have
$
 \sum_{s<t}\eta_s\leq 3t \lesssim\frac{\log N}{p_K}.
$
This time is within the horizon in
\Cref{prop:tail-v2}, because
\begin{align*}
 \frac{\log N}{p_K}
 &\lesssim(\log N)K^\alpha\lesssim(\log N)\left(\frac{d^2}{\log N}\right)^\alpha
 \leq d^{2\alpha+2}
\end{align*}
for sufficiently large d.

We have thus proved the following $K$-parameterized statement: simultaneously
for every integer $d\leq K\leq d^2/\log N$,
For every $t\lesssim\frac{\log N}{p_K},$ we have
\[
 L(\Wb_t^{\mathrm{GD}})
 \gtrsim Kp_K\log N.
\]
% Equivalently, if $\eps$ is at most a sufficiently small fixed multiple of
% $Kp_K\log N$, then $L(\Wb_t^{\mathrm{GD}})
%  >\eps$ and we get
% \[
%  \tau_{\mathrm{GD}}(\eps)
%  \gtrsim\frac{\log N}{p_K}.
% \]

% We now use the target $\eps$ to select $K$.  For an arbitrary
% $
%  C_{\mathrm{low}}
%  \frac{(\log N)^\alpha}{d^{2(\alpha-1)}}
%  \leq\eps\leq
%  c_{\mathrm{up}}
%  \frac{\log N}{d^{\alpha-1}},
% $
% choose an integer
% \[
%  K\eqsim
%  \left(\frac{\log N}{\eps}\right)^{1/(\alpha-1)},
% \]
% where the fixed implicit multiplicative constant is chosen sufficiently
% small. The target range, after increasing $C_{\mathrm{low}}$ and decreasing $c_{\mathrm{up}}$ to
% leave fixed-factor slack, therefore ensures
% $d\leq K\leq c d^2/\log N$.

% Moreover, 
% \begin{align*}
%  Kp_K\log N
%  &\eqsim(\log N)K^{1-\alpha}\eqsim\eps.
% \end{align*}
% Since $1-\alpha<0$, decreasing the fixed multiplicative constant in the
% choice of $K$ increases the ratio
% $(Kp_K\log N)/\eps$.  We may therefore choose it sufficiently small that
% $Kp_K\log N$ is a sufficiently large fixed multiple of $\eps$.  The
% $K$-parameterized lower bound then gives, for every
% $t\lesssim(\log N)/p_K$,
% \[
%  L(\Wb_t^{\mathrm{GD}})
%  \gtrsim Kp_K\log N
%  >\eps.
% \]
% None of these times can be the hitting time.  Hence
% \begin{align*}
%  \tau_{\mathrm{GD}}(\eps)
%  &\gtrsim\frac{\log N}{p_K}\eqsim(\log N)K^\alpha\eqsim(\log N)
%  \left(\frac{\log N}{\eps}\right)^{\alpha/(\alpha-1)}.
% \end{align*}
% Intersecting this pointwise schedule event with the common event in
% \Cref{thm:fw} and applying the Muon upper bound there proves
% $\tau_{\mathrm{GD}}(\eps)\gtrsim d\,\tau_{\mathrm{Muon}}(\eps)$.
% The probability statement remains pointwise in the fixed target-schedule
% pair because the leave-block trajectory is independent of the removed signs
% only for a schedule fixed before the sample.
We now eliminate $K$ in favor of the iteration number $t$.
Fix
\[
C_{\mathrm{low}}(\log N)d^\alpha
\leq t
\leq
c_{\mathrm{up}}
\frac{d^{2\alpha}}{(\log N)^{\alpha-1}},
\]
and choose an integer
\[
K\eqsim
\left(\frac{t}{\log N}\right)^{1/\alpha},
\]
where the fixed implicit multiplicative constant is chosen sufficiently
large. After increasing $C_{\mathrm{low}}$ and decreasing $c_{\mathrm{up}}$
to leave fixed-factor slack, this choice satisfies
\[
d\leq K\leq c\frac{d^2}{\log N}.
\]
Moreover, since $p_K\eqsim K^{-\alpha}$,
\[
\frac{\log N}{p_K}
\eqsim
(\log N)K^\alpha
\eqsim t.
\]
By choosing the implicit constant in $K$ sufficiently large, we can ensure
\[
t\lesssim \frac{\log N}{p_K},
\]
and hence the $K$-parameterized lower bound applies. Therefore
\begin{align*}
L(\Wb_t^{\mathrm{GD}})
&\gtrsim
Kp_K\log N\\
&\eqsim
(\log N)K^{1-\alpha}\\
&\eqsim
(\log N)
\left(\frac{t}{\log N}\right)^{-(\alpha-1)/\alpha}.
\end{align*}
Thus, with high probability, simultaneously for every
$
C_{\mathrm{low}}(\log N)d^\alpha
\leq t
\leq
c_{\mathrm{up}}
\frac{d^{2\alpha}}{(\log N)^{\alpha-1}},
$
we have
\[
L(\Wb_t^{\mathrm{GD}})
\gtrsim
(\log N)
\left(\frac{t}{\log N}\right)^{-(\alpha-1)/\alpha}.
\]

\end{proof}



% ============================================================
\section{Proof of \Cref{cor:gd-weight-decay}}
% ============================================================
\begin{proof}
For each $t$, define
\[
F_t(\Wb)
:=
L(\Wb)+\frac{\lambda_t}{2}\normF{\Wb}^2.
\]
On the regularity event of \Cref{lem:regularity-v2},
\Cref{lem:smoothness-v2} gives $\HF<2/3$, so $F_t$ is convex and
$(\HF+\lambda_t)$-smooth in Frobenius norm.  The assumed condition implies
\[
\eta_t
\leq
\frac{2}{\HF+\lambda_t},
\]
and hence the map
\[
\Wb\mapsto
\Wb-\eta_t\nabla F_t(\Wb)
=
(1-\eta_t\lambda_t)\Wb-\eta_t\nabla L(\Wb)
\]
is nonexpansive in Frobenius norm by \Cref{lem:cocoercive-v2}.

For a block $B$, the corresponding leave-block update can be written as
the same nonexpansive map plus
$\eta_t\nabla L_B$.  Therefore the leave-block stability bound in
\Cref{lem:block-v2} holds unchanged.  Moreover, the assumption above
implies $\eta_t\leq3$, so all subsequent tail-preservation and
$K$-parameterized lower-bound arguments apply verbatim.  The claimed
learning-curve and hitting-time bounds then follow exactly as in
\Cref{thm:gd}.
\end{proof}

% ============================================================
\section{Proof of \Cref{thm:capped}}
% ============================================================

\begin{proof}[Proof of \Cref{thm:capped}]
Recall from the proof of \Cref{thm:Muon-upper bound} that $M:=(\log N)d^{\alpha-1},$
and that the radius schedule is
\[
R_t
=
C_\alpha(\log N)
\max\left\{
1,
\left(\frac{t+1}{M}\right)^{1/(2\alpha)}
\right\}.
\]
The only change from \Cref{thm:Muon-upper bound} is the mixing schedule:
\[
\gamma_t=\frac{2}{t+2+a},
\qquad
a=
\max\left\{
0,\frac{2C_\alpha\log N}{\bar\eta}-2
\right\}.
\]
In particular, \(a\ge0\), so \(0<\gamma_t\le1\).

We first prove the deterministic effective-learning-rate bound. By construction
\begin{equation}
\frac{2C_\alpha\log N}{a+2}\le\bar\eta.
\label{eq:capped-a-choice}
\end{equation}
Moreover,
\[
a+2
\le
2+\frac{2C_\alpha\log N}{\bar\eta}
=
O_{\bar\eta}(\log N),
\]
whereas \(M=(\log N)d^{\alpha-1}\).  Since \(\alpha>1\), for all
sufficiently large \(d\),
\begin{equation}
M\ge a+2.
\label{eq:capped-M-a}
\end{equation}

We now show that \(\gamma_tR_t\le\bar\eta\) for every \(t\ge0\).
If \(t+1\le M\), then \(R_t=C_\alpha\log N\), and therefore
\[
\gamma_tR_t
=
\frac{2C_\alpha\log N}{t+2+a}
\le
\frac{2C_\alpha\log N}{a+2}
\le
\bar\eta.
\]
If \(t+1>M\), then, since \(1/(2\alpha)<1\),
\[
\left(\frac{t+1}{M}\right)^{1/(2\alpha)}
\le
\frac{t+1}{M}.
\]
Hence, using \eqref{eq:capped-M-a},
\begin{align*}
\gamma_tR_t
&=
\frac{2C_\alpha\log N}{t+2+a}
\left(\frac{t+1}{M}\right)^{1/(2\alpha)}\\
&\le
\frac{2C_\alpha\log N}{a+2}
\frac{t+1}{t+2+a}\\
&\le
\frac{2C_\alpha\log N}{a+2}
\le
\bar\eta.
\end{align*}
Thus, for all sufficiently large \(d\),
\begin{equation}
\gamma_tR_t\le\bar\eta
\qquad\text{for every }t\ge0.
\label{eq:capped-effective}
\end{equation}

We next control the actual displacement of the Muon iterates.  Since
\(R_t\) is nondecreasing, the same feasibility induction used in
\Cref{lem:Muon-changing-radius} applies.  Starting from
\(\Wb_0^{\mathrm{Muon}}=\mathbf0\), if
\(\normop{\Wb_t^{\mathrm{Muon}}}\le R_t\), then the operator-ball oracle
\[
\Sb_t
:=
-R_t\polar\!\left(\nabla L(\Wb_t^{\mathrm{Muon}})\right)
\]
satisfies \(\normop{\Sb_t}\le R_t\), and therefore
\begin{align*}
\normop{\Wb_{t+1}^{\mathrm{Muon}}}
&=
\normop{
(1-\gamma_t)\Wb_t^{\mathrm{Muon}}
+\gamma_t\Sb_t
}\\
&\leq
(1-\gamma_t)\normop{\Wb_t^{\mathrm{Muon}}}
+\gamma_t\normop{\Sb_t}\\
&\leq R_t
\leq R_{t+1}.
\end{align*}
Hence
\[
\normop{\Wb_t^{\mathrm{Muon}}}\le R_t
\qquad
\text{for every }t\ge0.
\]
Consequently,
\begin{align*}
\normop{
\Wb_{t+1}^{\mathrm{Muon}}
-\Wb_t^{\mathrm{Muon}}
}
&=
\gamma_t
\normop{
\Sb_t-\Wb_t^{\mathrm{Muon}}
}\\
&\leq
\gamma_t
\left(
\normop{\Sb_t}
+\normop{\Wb_t^{\mathrm{Muon}}}
\right)\\
&\leq
2\gamma_tR_t\\
&\leq2\bar\eta,
\end{align*}
where the last step uses \eqref{eq:capped-effective}.

It remains to prove the loss bound.  We work on the same event used in the
proof of \Cref{thm:Muon-upper bound}: the intersection of the event in
\Cref{prop:prefix-v2} and the regularity event in
\Cref{lem:regularity-v2}.  This event has probability at least
\(1-O(d^{-q})\), and on it, $\Hop\leq\frac{1221}{2000}.$
No new probabilistic event is needed below.

Fix \(t\) in the iteration range stated in the theorem.  In particular,
\(t\) lies in the range of \Cref{thm:Muon-upper bound}, so, after the same
choice of endpoint constants as there,
\[
t\geq M.
\]
For every \(0\le s<t\), use exactly the same prefix comparator
\(\widehat\Wb_{K_s}\) as in \Cref{lem:Muon-time-comparators}.  Thus
\begin{align}
\normop{\widehat\Wb_{K_s}}
&\leq R_s,
\label{eq:capped-comparator-feasible}\\
L(\widehat\Wb_{K_s})
&\leq
A_{\alpha,q}(\log N)d^{1-\alpha}
\max\left\{
1,\frac{s+1}{M}
\right\}^{-(\alpha-1)/\alpha}.
\label{eq:capped-comparator-loss}
\end{align}

The exact linear minimization oracle, comparator feasibility
\eqref{eq:capped-comparator-feasible}, convexity, and operator smoothness
give exactly as in the proof of \Cref{lem:Muon-changing-radius}
\begin{align*}
L(\Wb_{s+1}^{\mathrm{Muon}})
&\leq
(1-\gamma_s)L(\Wb_s^{\mathrm{Muon}})
+\gamma_sL(\widehat\Wb_{K_s})
+2\Hop\gamma_s^2R_s^2\\
&=
\frac{s+a}{s+a+2}L(\Wb_s^{\mathrm{Muon}})
+\frac{2}{s+a+2}L(\widehat\Wb_{K_s})
+\frac{8\Hop}{(s+a+2)^2}R_s^2.
\end{align*}
Multiplying by \((s+a+1)(s+a+2)\) yields
\begin{align*}
&(s+a+1)(s+a+2)L(\Wb_{s+1}^{\mathrm{Muon}})\\
&\qquad\leq
(s+a)(s+a+1)L(\Wb_s^{\mathrm{Muon}})
+2(s+a+1)L(\widehat\Wb_{K_s})+
8\Hop\frac{s+a+1}{s+a+2}R_s^2.
\end{align*}
Summing from \(s=0\) to \(s=t-1\) telescopes and gives
\begin{align}
&(t+a)(t+a+1)L(\Wb_t^{\mathrm{Muon}})\nonumber\\
&\qquad\leq
a(a+1)L(\Wb_0^{\mathrm{Muon}})
+2\sum_{s=0}^{t-1}
(s+a+1)L(\widehat\Wb_{K_s})
+8\Hop\sum_{s=0}^{t-1}\frac{s+a+1}{s+a+2}R_s^2.\label{eq:capped-master}
\end{align}
Since \(L(\Wb_0^{\mathrm{Muon}})=\log N\) and
\((t+a)(t+a+1)\ge t^2\), it remains to bound the three terms on the
right-hand side.

We begin with the comparator contribution.  Write $p:=\frac{\alpha-1}{\alpha}\in(0,1).$
Using \eqref{eq:capped-comparator-loss},
\begin{align*}
&\sum_{s=0}^{t-1}
(s+a+1)L(\widehat\Wb_{K_s})\\
&\qquad\leq
A_{\alpha,q}(\log N)d^{1-\alpha}
\sum_{s=0}^{t-1}
(s+a+1)
\max\left\{
1,\frac{s+1}{M}
\right\}^{-p}\\
&\qquad=
A_{\alpha,q}(\log N)d^{1-\alpha}
\left[
\sum_{s=0}^{t-1}
(s+1)
\max\left\{
1,\frac{s+1}{M}
\right\}^{-p}+
a\sum_{s=0}^{t-1}
\max\left\{
1,\frac{s+1}{M}
\right\}^{-p}
\right].
\end{align*}
The first sum is exactly the first scalar sum in
\Cref{lem:Muon-scalar-sums}:
\[
\sum_{s=0}^{t-1}
(s+1)
\max\left\{
1,\frac{s+1}{M}
\right\}^{-p}
\leq
M^p t^{1+1/\alpha}.
\]
For the second sum, setting \(m=s+1\) and using
\[
\max\left\{1,\frac mM\right\}^{-p}
\leq M^pm^{-p},
\]
we obtain
\begin{align*}
\sum_{s=0}^{t-1}
\max\left\{
1,\frac{s+1}{M}
\right\}^{-p}
&\leq
M^p\sum_{m=1}^tm^{-p}\leq
\alpha M^pt^{1/\alpha}.
\end{align*}
Therefore
\begin{align}
\sum_{s=0}^{t-1}
(s+a+1)L(\widehat\Wb_{K_s})
&\leq
A_{\alpha,q}(\log N)d^{1-\alpha}
M^pt^{1+1/\alpha}
\left(
1+\alpha\frac at
\right).
\label{eq:capped-comparator-sum}
\end{align}
Since
\[
a=O_{\bar\eta}(\log N)
\qquad\text{and}\qquad
t\geq M=(\log N)d^{\alpha-1},
\]
we have
\begin{equation}
\frac at
\lesssim_{\bar\eta}
d^{-(\alpha-1)}
=o(1).
\label{eq:capped-a-over-t}
\end{equation}
Hence
\begin{equation}
\sum_{s=0}^{t-1}
(s+a+1)L(\widehat\Wb_{K_s})
\lesssim_{\bar\eta}
(\log N)d^{1-\alpha}
M^pt^{1+1/\alpha}.
\label{eq:capped-comparator-final}
\end{equation}

For the curvature contribution, the shift causes no loss in the estimate,
because
\[
\frac{s+a+1}{s+a+2}\leq1.
\]
Using the same radius identity and the same second scalar sum as in the
proof of \Cref{thm:Muon-upper bound},
\begin{align}
\sum_{s=0}^{t-1}
\frac{s+a+1}{s+a+2}R_s^2
&\leq
C_\alpha^2(\log N)^2
\sum_{s=0}^{t-1}
\max\left\{
1,\frac{s+1}{M}
\right\}^{1/\alpha}
\notag\\
&\leq
C_\alpha^2(\log N)^2
M^{-1/\alpha}t^{1+1/\alpha}.
\label{eq:capped-curvature}
\end{align}

Finally, we control the new initial term created by the shifted
telescoping.  Since \(a=O_{\bar\eta}(\log N)\),
\[
\frac{a(a+1)\log N}{t^2}
\lesssim_{\bar\eta}
\frac{(\log N)^3}{t^2}.
\]
We compare this with the scale of the comparator term.  Since \(t\ge M\)
and \(p-2<0\),
\begin{align*}
\frac{
a(a+1)(\log N)t^{-2}
}{
(\log N)d^{1-\alpha}M^pt^{-p}
} &=
a(a+1)d^{\alpha-1}M^{-p}t^{p-2}\\
&\leq
a(a+1)d^{\alpha-1}M^{-2}\\
&=
\frac{a(a+1)}
{(\log N)^2d^{\alpha-1}}
\lesssim_{\bar\eta}
d^{-(\alpha-1)}
=o(1).
\end{align*}
Thus
\begin{equation}
\frac{a(a+1)\log N}{t^2}
\lesssim_{\bar\eta}
(\log N)d^{1-\alpha}M^pt^{-p}.
\label{eq:capped-initial-term}
\end{equation}

Substituting
\eqref{eq:capped-comparator-final},
\eqref{eq:capped-curvature}, and
\eqref{eq:capped-initial-term}
into \eqref{eq:capped-master}, and using
\((t+a)(t+a+1)\ge t^2\), gives
\begin{align*}
L(\Wb_t^{\mathrm{Muon}})
&\lesssim_{\bar\eta}
(\log N)d^{1-\alpha}
M^pt^{-p}\\
&\qquad+
(\log N)^2
M^{-1/\alpha}t^{-p}.
\end{align*}
Exactly as in the final simplification in the proof of
\Cref{thm:Muon-upper bound}, the two coefficients have the same scale:
\begin{align*}
(\log N)^2M^{-1/\alpha}
&=
(\log N)d^{1-\alpha}M^{(\alpha-1)/\alpha}.
\end{align*}
Therefore
\begin{align*}
L(\Wb_t^{\mathrm{Muon}})
&\lesssim_{\bar\eta}
(\log N)d^{1-\alpha}
M^{(\alpha-1)/\alpha}
t^{-(\alpha-1)/\alpha}=
(\log N)\left(\frac{dt}{\log N}\right)^{-(\alpha-1)/\alpha}.
\end{align*}

The comparator event used above is the same single event from
\Cref{prop:prefix-v2} used in the proof of
\Cref{thm:Muon-upper bound}, and all remaining steps are deterministic.
Hence the loss bound holds simultaneously for every \(t\) in the stated
iteration range on an event of probability at least \(1-O(d^{-q})\).
This completes the proof.
\end{proof}

% ============================================================
% \section{Proof of \Cref{thm:capped}}
% % ============================================================

% \begin{proof}[Proof of \Cref{thm:capped}]
% Work on the common event of probability $1-O(d^{-q})$ from
% \Cref{thm:fw}.  Only the mixing schedule changes; the operator ball and its
% polar oracle are the same.  Since $a\geq0$, one has
% $0<\gamma_t\leq1$.  Moreover, $t\mapsto\gamma_t$ is decreasing and the
% chosen $a$ satisfies
% $2+a\geq2R_{\mathrm{op}}(\eps/2)/\bar\eta$,
% \begin{align*}
%  R_{\mathrm{op}}(\eps/2)\gamma_t
%  &=\frac{2R_{\mathrm{op}}(\eps/2)}{t+2+a}\leq\frac{2R_{\mathrm{op}}(\eps/2)}{2+a}\leq\bar\eta.
% \end{align*}
% Every oracle point satisfies
% $\normop{\Sb_t^{\mathrm{op}}}\leq R_{\mathrm{op}}(\eps/2)$.  Starting from
% $\Wb_0^{\mathrm{op}}=\mathbf0$, feasibility follows by induction: if the
% $t$th iterate is feasible, then
% \begin{align*}
%  \normop{\Wb_{t+1}^{\mathrm{op}}}
%  &=\normop{
%   (1-\gamma_t)\Wb_t^{\mathrm{op}}+\gamma_t\Sb_t^{\mathrm{op}}
%  }\\
%  &\leq
%  (1-\gamma_t)\normop{\Wb_t^{\mathrm{op}}}
%  +\gamma_t\normop{\Sb_t^{\mathrm{op}}}\\
%  &\leq
%  (1-\gamma_t)R_{\mathrm{op}}(\eps/2)
%  +\gamma_tR_{\mathrm{op}}(\eps/2)\\
%  &=R_{\mathrm{op}}(\eps/2).
% \end{align*}
% Thus both the iterate and oracle point lie in the operator ball.  Therefore
% \begin{align*}
%  \normop{\Wb_{t+1}^{\mathrm{op}}-\Wb_t^{\mathrm{op}}}
%  &=\gamma_t\normop{\Sb_t^{\mathrm{op}}-\Wb_t^{\mathrm{op}}}\\
%  &\leq\gamma_t
%  \left(\normop{\Sb_t^{\mathrm{op}}}+\normop{\Wb_t^{\mathrm{op}}}\right)\\
%  &\leq2R_{\mathrm{op}}(\eps/2)\gamma_t\leq2\bar\eta.
% \end{align*}

% Let $\Wb_R^*$ minimize $L$ on this ball and put
% $h_t:=L(\Wb_t^{\mathrm{op}})-L(\Wb_R^*)$.  The ball contains a point of
% loss at most $\eps/2$, so $L(\Wb_R^*)\leq\eps/2$.  Repeating the oracle and
% convexity argument with every line displayed gives
% \begin{align*}
%  &\left\langle
%   \nabla L(\Wb_t^{\mathrm{op}}),
%   \Sb_t^{\mathrm{op}}-\Wb_t^{\mathrm{op}}
%  \right\rangle_{\mathrm F}\\
%  &\quad=
%  \min_{\normop{\Sb}\leq R_{\mathrm{op}}(\eps/2)}
%  \left\langle\nabla L(\Wb_t^{\mathrm{op}}),\Sb\right\rangle_{\mathrm F}
%  -\left\langle
%   \nabla L(\Wb_t^{\mathrm{op}}),\Wb_t^{\mathrm{op}}
%  \right\rangle_{\mathrm F}\\
%  &\quad\leq
%  \left\langle
%   \nabla L(\Wb_t^{\mathrm{op}}),\Wb_R^*-\Wb_t^{\mathrm{op}}
%  \right\rangle_{\mathrm F}\\
%  &\quad\leq L(\Wb_R^*)-L(\Wb_t^{\mathrm{op}})=-h_t.
% \end{align*}
% The last inequality is exactly the convexity inequality
% \[
%  L(\Wb_R^*)\geq L(\Wb_t^{\mathrm{op}})
%  +\left\langle
%   \nabla L(\Wb_t^{\mathrm{op}}),\Wb_R^*-\Wb_t^{\mathrm{op}}
%  \right\rangle_{\mathrm F}.
% \]

% The displacement identity and the ball diameter give
% \begin{align*}
%  \Wb_{t+1}^{\mathrm{op}}-\Wb_t^{\mathrm{op}}
%  &=\gamma_t(\Sb_t^{\mathrm{op}}-\Wb_t^{\mathrm{op}}),\\
%  \normop{\Sb_t^{\mathrm{op}}-\Wb_t^{\mathrm{op}}}
%  &\leq2R_{\mathrm{op}}(\eps/2).
% \end{align*}
% Substituting these facts into operator smoothness yields
% \begin{align*}
%  h_{t+1}
%  &\leq h_t
%  +\gamma_t\left\langle
%   \nabla L(\Wb_t^{\mathrm{op}}),
%   \Sb_t^{\mathrm{op}}-\Wb_t^{\mathrm{op}}
%  \right\rangle_{\mathrm F}+
%  \frac{\Hop\gamma_t^2}{2}
%  \normop{\Sb_t^{\mathrm{op}}-\Wb_t^{\mathrm{op}}}^2\\
%  &\leq(1-\gamma_t)h_t
%  +2\Hop R_{\mathrm{op}}(\eps/2)^2\gamma_t^2\\
%  &=\frac{t+a}{t+a+2}h_t
%  +\frac{8\Hop R_{\mathrm{op}}(\eps/2)^2}{(t+a+2)^2}.
% \end{align*}

% Multiply the recurrence at time $s$ by $(s+a+1)(s+a+2)$.  Summing the
% resulting differences explicitly gives
% \begin{align*}
%  &(t+a)(t+a+1)h_t-a(a+1)h_0\\
%  &\quad=
%  \sum_{s=0}^{t-1}
%  \Bigl[
%   (s+a+1)(s+a+2)h_{s+1}
%   -(s+a)(s+a+1)h_s
%  \Bigr]\\
%  &\quad\leq
%  8\Hop R_{\mathrm{op}}(\eps/2)^2
%  \sum_{s=0}^{t-1}\frac{s+a+1}{s+a+2}\\
%  &\quad\leq8\Hop R_{\mathrm{op}}(\eps/2)^2t.
% \end{align*}
% Because $h_0=L(\mathbf0)-L(\Wb_R^*)\leq\log N$, for every $t\geq1$ division
% by $(t+a)(t+a+1)\geq t^2$ gives
% \begin{equation}\label{eq:h_t}
%  h_t
%  \leq
%  \frac{a(a+1)\log N}{t^2}
%  +\frac{8\Hop R_{\mathrm{op}}(\eps/2)^2}{t}.
% \end{equation}

% The choice of $a$ gives
% $a+1\le 2R_{\mathrm{op}}(\eps/2)/\bar \eta$. Therefore we have
% $
%  a(a+1)\le 4R_{\mathrm{op}}(\eps/2)^2/(\bar \eta)^2.
% $
% Putting this back into Eq.~\ref{eq:h_t} and get,
% \begin{equation}\nonumber
%  h_t
%  \leq
%  \frac{4R_{\mathrm{op}}(\eps/2)^2\log N}{\bar \eta^2t^2}
%  +\frac{8\Hop R_{\mathrm{op}}(\eps/2)^2}{t}.
% \end{equation}
% To make the first term at most $\eps/4$, it suffices to require
% \[
%  \frac{4R_{\mathrm{op}}(\eps/2)^2\log N}
%  {\bar\eta^2 t^2}
%  \leq \frac{\eps}{4},
% \]
% or equivalently,
% \[
%  t\geq
%  t_1:=
%  \frac{4R_{\mathrm{op}}(\eps/2)}{\bar\eta}
%  \sqrt{\frac{\log N}{\eps}}.
% \]
% Similarly, to make the second term at most $\eps/4$, it suffices to require
% \[
%  \frac{8\Hop R_{\mathrm{op}}(\eps/2)^2}{t}
%  \leq\frac{\eps}{4},
% \]
% or equivalently,
% \[
%  t\geq
%  t_2:=
%  \frac{32\Hop R_{\mathrm{op}}(\eps/2)^2}{\eps}.
% \]

% We now compare these two time scales.  Their ratio is
% \begin{align*}
%  \frac{t_1}{t_2}
%  &=
%  \frac{
%  \frac{4R_{\mathrm{op}}(\eps/2)}{\bar\eta}
%  \sqrt{\frac{\log N}{\eps}}
%  }{
%  \frac{32\Hop R_{\mathrm{op}}(\eps/2)^2}{\eps}
%  }
%  =
%  \frac{1}{8\bar\eta\Hop}
%  \frac{\sqrt{\eps\log N}}
%  {R_{\mathrm{op}}(\eps/2)}.
% \end{align*}
% The sharp operator radius formula and the upper endpoint of the target range give \begin{align*} 
% R_{\mathrm{op}}(\eps/2)^2 &\eqsim \frac{(\log N)^2}{d} \left(\frac{\log N}{\eps/2}\right)^{1/(\alpha-1)}\\ &\gtrsim \frac{(\log N)^2}{d}\,d =(\log N)^2. \end{align*} 
% Thus $R_{\mathrm{op}}(\eps/2)\gtrsim\log N$
% and $\Hop\eqsim1$.  Hence
% \[
%  \frac{t_1}{t_2}
%  \lesssim_{\bar\eta}
%  \sqrt{\frac{\eps}{\log N}}=o(1).
% \]

% In particular, for all sufficiently large $d$, taking $t\geq t_2$
% automatically ensures $t\geq t_1$ as well.  Consequently both terms are at
% most $\eps/4$, and hence $h_t\leq\frac{\eps}{2}.$ 
% \[
%  L(\Wb_t^{\mathrm{op}})
%  =L(\Wb_R^*)+h_t
%  \leq\frac\eps2+\frac\eps2=\eps.
% \]
% Consequently,
% \[
%  \tau_{\mathrm{Muon}}(\eps;\gamma)
%  \lesssim t_2\eqsim \frac{\Hop R_{\mathrm{op}}(\eps/2)^2}{\eps} \lesssim \frac{\log N}{d}
%  \left(\frac{\log N}{\eps}\right)^{\alpha/(\alpha-1)}.
% \]
% This proves the theorem.
% \end{proof}

\begin{thebibliography}{9}

\bibitem{vershynin}
R.~Vershynin,
\emph{High-Dimensional Probability: An Introduction with Applications in
Data Science}, Cambridge University Press, 2018.

\bibitem{hoeffding}
W.~Hoeffding,
Probability inequalities for sums of bounded random variables,
\emph{Journal of the American Statistical Association} 58 (1963), 13--30.

\end{thebibliography}

\end{document}
